% ============================================================
%  Body — The Venus Ring of Voynich f57v
%  Outline / skeleton with section headers, one-line summaries,
%  and content notes (in LaTeX comments) for each subsection.
%  Sections to be drafted in order after author approval.
% ============================================================

% ============================================================
\section{Introduction}
\label{sec:intro}

The Voynich Manuscript (Beinecke Rare Book and Manuscript Library MS~408)
is a hand-drawn codex of approximately 234 surviving folios, radiocarbon-dated
to 1404--1438~CE, written in an unknown script and accompanied by an
extensive program of botanical, cosmological, balneological, and
pharmaceutical illustrations. It has resisted decipherment for over a
century. Most published proposals to read it as a substitution cipher of
some natural language have failed under quantitative scrutiny, and most
proposals to identify a specific textual source for its content have
failed under documentary scrutiny. The most consequential negative result
in the field is that no internally consistent reading has ever been
produced for any continuous passage of the manuscript longer than a few
words.

This paper is the third in a series of computational studies of the
Voynich Manuscript by the present author. The first two
papers~\citep{jimenez2026negative,jimenez2026twotier} characterised the
manuscript's global text structure: the first established that nineteen
substitution-based decipherment models against five Semitic languages
all converge on the same structural ceiling without producing readable
text, and the second established that the manuscript's positional
structure is consistent with a designed two-tier combinatorial code
rather than with the surface form of an underlying natural-language
text. Both papers characterised global properties of the manuscript;
neither identified a specific historical or textual tradition to which
the manuscript belongs. The interpretive question---\emph{what kind of
medieval document is this, and where does it sit in the documented
textual landscape of the early fifteenth century?}---was left open.

The present paper takes a different methodological approach to that
question. Rather than continuing to analyse global structure, it
identifies and quantitatively characterises a single \emph{local
anomaly} in the manuscript and traces it to a documented historical
tradition. Specifically, we report that folio f57v---a cosmological
circular diagram showing four nymphs around a central solar
symbol~$\odot$---contains a private subalphabet of five glyphs that
appear nowhere else in the entire 37,087-word STA1 transliteration of
the manuscript, plus three further glyphs concentrated on f57v at
enrichment factors of $75\times$ to $302\times$. We show that the
concentration is statistically incompatible with chance by tens of
orders of magnitude, that it is geographically confined to the four
concentric text rings of f57v (most densely to line f57v.3), and that
it is absent from every other folio of the manuscript's
cosmological/astronomical block (the foldout pages of ff.~67--70). We
then demonstrate that this localized inscription coexists on the same
folio with the manuscript's general categorial labeling system---which
is fully present in f57v's central nymph labels---without contaminating
it.

The two coexisting layers correspond, we argue, to a documented
medieval textual category: the descriptive astrological text plus the
reserved talismanic inscription in non-standard characters
(\textit{karacteres}) prescribed by the Latin Picatrix tradition for
the manufacture of planetary images. We support this interpretation by
direct visual comparison with the only extant fully illustrated copy of
\textit{Picatrix Latinus}, Cracow Jagiellonian Library MS~793 (Piotr
Gaszowiec, 1471--1474, copied from a reconstructed Bohemian prototype
c.~1390--1410), and by direct verification against the printed
critical edition of the Latin Picatrix text by \citet{pingree1986}.

The paper makes two main claims. The \textbf{empirical claim}
(\S\S\ref{sec:subalphabet}--\ref{sec:two_phenomena}) is that the
Voynich Manuscript contains two textually distinct phenomena that
coexist in the same codex without contaminating each other:

\begin{itemize}[itemsep=0pt,topsep=2pt]
\item \textbf{Phenomenon~A}, a general categorial labeling system
operating throughout most of the manuscript and validated for the
zodiac block, the f57v central nymph labels, the f70r1 sun-ray
labels, and body text in the herbal/pharmaceutical/recipe sections;
\item \textbf{Phenomenon~B}, a localized enumerative inscription
composed of a private subalphabet of eight glyphs, confined to the
four concentric text rings of folio f57v and most densely to line
f57v.3, and absent from every other folio of the manuscript.
\end{itemize}

The empirical claim is independent of any interpretive frame and is
defended on purely internal evidence from the STA1 transliteration of
the manuscript and from the IVTFF locus-tag system.

The \textbf{interpretive claim} (\S\S\ref{sec:verification}--\ref{sec:discussion})
is that Phenomenon~B is plausibly an instance of the medieval
astrological-talismanic \textit{karacteres} tradition codified by the
Latin Picatrix, and that the Voynich Manuscript should therefore be
read as an \emph{independent witness} of that tradition---not as a
copy or translation of any specific Picatrix manuscript, but as a
sibling production of the same documented textual milieu, datable
within the Voynich's C14 window to the same Toledo--Prague--Cracow
geographic and chronological context in which the Picatrix tradition
is known to have been actively practised. The interpretive claim is
explicitly weaker than ``decipherment'': it does not assert
character-level correspondences, textual lineage, or an underlying
natural language. It asserts \emph{shared membership in a documented
tradition}, with the structural anomaly of f57v as the empirical
anchor.

\subsection{Strategic shift: from global structure to local anomaly}
\label{sec:strategic_shift}

The methodological strategy of this paper is the opposite of the two
companion studies. Where papers~1 and~2 sought general structural
patterns across the entire 224-folio corpus and asked whether those
patterns were compatible with various global hypotheses, the present
paper deliberately seeks the \emph{single most distinctive point} in
the manuscript and asks what it points to.

The reason for the shift is methodological. Papers~1 and~2 produced
strong negative and structural results, but these results
underdetermine the historical question. A two-tier positional code
is consistent with many different medieval textual genres
(encyclopaedic indices, lapidaries, astronomical tables, alchemical
recipes, talismanic manuals), and global structural arguments cannot
discriminate between them. Local anomalies, by contrast, are
discriminative precisely because they are unusual: a folio that
behaves textually unlike any other folio in the manuscript is, by
construction, a folio whose textual signature can be matched against
specific categories rather than against general statistics.

Folio f57v turns out to be such an anomaly. As \S\ref{sec:ranking}
will document, it is the rank-1 folio in the entire manuscript by
proportion of isolated single-character ``words'' ($63.9\%$, against
a corpus median of $2.3\%$), and it is the only folio whose dominant
textual behaviour is enumerative rather than continuous. This kind of
anomaly is the right kind of empirical handle for a historical
identification: there is exactly one folio with this signature, and
the question is whether any documented medieval textual category
prescribes \emph{exactly that} signature in \emph{exactly that}
iconographic context. The remainder of the paper develops the answer.

\subsection{Notation, corpus, and reproducibility}
\label{sec:notation}

Throughout this paper we work from Ren\'e Zandbergen's STA1~2.0 super
transliteration of the Voynich Manuscript~\citep{zandbergen_sta1}, a
fine-grained transcription of the manuscript's atomic glyph forms
into ASCII codes. The corpus contains $37{,}087$ word tokens,
$157{,}254$ atom occurrences, and $166$ distinct atom types
distributed over $224$ folios. STA1 makes finer distinctions than the
more widely cited EVA transliteration, and the present paper depends
on those distinctions in ways discussed in \S\ref{sec:eva_caveat}.

Locus tags follow the IVTFF format
specification~\citep{zandbergen_ivtff}, in which each line of the
manuscript is identified by a folio name, a line number, and a
two-character tag indicating the page element (\texttt{@P0} for
ordinary paragraph text, \texttt{@L0} for centred labels, \texttt{@Ro}
for rotated radial labels, \texttt{@Cc}/\texttt{+Cc} for concentric
circular text, \texttt{@Ri} for radial text inside a diagram, etc.).
This tag system is essential for the analysis of \S\ref{sec:intrafolio}
because the f57v anomaly is page-element-specific, not folio-uniform:
it lives in the \texttt{@Cc}/\texttt{+Cc} concentric ring lines and
not in any \texttt{@L0}/\texttt{@Ro} label of the same folio.

All quantitative results in this paper are reproducible from the
project's public GitHub repository, which contains the principal
companion notebook (\filename{74_f57v_private_subalphabet.ipynb}),
the structured JSON outputs used as the source of every numerical
table and figure, and the standalone polar rectification tool
(\filename{tools/circular_text_unwrap.ipynb}) of
\S\ref{sec:vc4_alphabet}. The full inventory of artefacts and
external sources is collected in \S\ref{sec:data_availability}.


% ============================================================
\section[Theoretical framework: \textit{cifra ineffabile} and the
talismanic tradition]{Theoretical framework: \textit{cifra ineffabile}\\
and the medieval astrological-talismanic tradition}
\label{sec:framework}

This section establishes the theoretical category to which the empirical
findings of the rest of the paper will be assigned. The category is the
medieval astrological-talismanic tradition of \textit{cifras ineffabiles}
or ``ineffable ciphers'': closed inventories of non-alphabetic graphic
characters whose function is not phonetic transcription but operative
encoding within a categorial system. This is the kind of textual object
that Phenomenon~B (\S\ref{sec:phenomenonB}) most closely resembles, and
it is the kind of object that the Picatrix tradition explicitly
prescribes. Three things need to be in place before that argument can
be made: a clear distinction between substitution-based ciphers and
categorial pasigraphies (\S\ref{sec:cat_encoding}); a brief situation
of the Picatrix corpus within its documented transmission history
(\S\ref{sec:picatrix_corpus}); and an explicit characterisation of the
\textit{karacteres} prescription within the Picatrix text itself
(\S\ref{sec:karacteres_prescription}).

\subsection{From substitution ciphers to categorial encoding}
\label{sec:cat_encoding}

The default assumption in nearly a century of Voynich scholarship has
been that the manuscript is a \emph{substitution cipher}: a writing
system in which each Voynich glyph stands, by some fixed but unknown
rule, for a letter, syllable, or phoneme of an underlying natural
language. The first companion paper of this series~\citep{jimenez2026negative}
tested nineteen different substitution-based decipherment models against
five Semitic languages (Hebrew, Aramaic, Arabic, Judeo-Arabic, Geez)
and found that all of them converge on the same structural ceiling
without producing readable text. The second companion
paper~\citep{jimenez2026twotier} established that the Voynich's
positional structure is a designed two-tier categorial code rather
than the surface form of an underlying natural-language text.

The cumulative force of these results is that the Voynich is not a
substitution cipher of any tested natural language. The alternative
theoretical category we adopt in the present paper is that of
\emph{categorial} or \emph{pasigraphic} encoding: a writing system in
which the position of a sign within a structured slot carries the
semantic weight, rather than the sign itself standing for a specific
phoneme. In a categorial encoding, the ``alphabet'' is not a phonetic
inventory but a small set of conventional marks distributed across
positionally meaningful slots; reading the system requires knowledge
of the slot structure rather than knowledge of an underlying language.

The most explicit early-modern formulation of this alternative is
Juan Caramuel y Lobkowitz's \textit{Apparatus philosophicus} (1665),
which introduces the term \textit{cifra ineffabilis} (``ineffable
cipher'') for writing systems whose elements do not correspond to
spoken sounds but to categories within a formal system. Caramuel
develops the idea in dialogue with Lullian combinatorial logic, with
Aristotelian categories, and with the Kabbalistic letter-mysticism of
the late Renaissance. The relevant point for the present paper is
that Caramuel \emph{recovers} a medieval tradition rather than
inventing one: his \textit{cifras ineffabiles} are an early-modern
codification of practices that the medieval astrological-talismanic
manuscripts (Picatrix among them) had already been documenting for
several centuries. The recent academic recovery of Caramuel's
linguistic framework by \citet{martinezgavilan2025} establishes the
philological lineage; the project's own theoretical companion document
on Caramuel's \textit{Apparatus philosophicus} extends the connection
to the Voynich case.

The Voynich Manuscript's positional code structure, as documented in
the second companion paper~\citep{jimenez2026twotier}, places it
\emph{within} the category of categorial pasigraphies rather than
within the category of substitution ciphers. The present paper extends
this picture: it identifies, on a single folio of the manuscript,
a localized inscription (Phenomenon~B, \S\ref{sec:phenomenonB}) of
the type that Caramuel's framework would call a \textit{cifra
ineffabilis} in the strict sense---a closed inventory of
non-phonetic characters in cyclic enumerative deployment. The general
labeling layer of the manuscript (Phenomenon~A) is the broader
categorial encoding; the f57v ring inscription is its strict
talismanic instance.

\subsection{The Picatrix corpus and its medieval transmission}
\label{sec:picatrix_corpus}

The Picatrix corpus is the principal medieval witness of the categorial
talismanic tradition we have just described. It is also one of the
best documented, and its transmission history is unusually well
attested. We summarise the key facts here; the historical-geographic
detail relevant to the Voynich's C14 window has already been laid out
in \S\ref{sec:cracow_milieu} and \S\ref{sec:transmission} and is not
repeated.

\noindent\textbf{The Arabic original.} The
\textit{Gh\=ayat al-\d{H}ak\={\i}m} (``The Aim of the Sage'') is an
Arabic compilation on astral magic, astrology, and operative talismanic
practice. It was traditionally attributed to the Andalusi mathematician
and astronomer \textit{Maslama b.~Ahmad al-Majr\={\i}\d{t}\={\i}}
(d.~1007), but \citet{fierro1996} reattributes it to the earlier
Cordoban polymath \textit{Maslama b.~Q\=asim al-Qur\d{t}ub\={\i}}
(d.~353/964), the same author as the \textit{Rutbat al-\d{H}ak\={\i}m}.
The reattribution places the composition of the original text squarely
in tenth-century al-Andalus, within the Batinist--Ismaili intellectual
network of Maslama al-Qur\d{t}ub\={\i}'s circle.

\noindent\textbf{The Castilian and Latin translations.} Under the
patronage of King Alfonso~X of Castile (Alfonso ``the Wise''), the
\textit{Gh\=ayat al-\d{H}ak\={\i}m} was translated into Castilian
c.~1256--1258 by an anonymous translator working in the orbit of the
Toledo translation school. A Latin translation was produced shortly
afterwards, almost certainly from the Castilian intermediate
(c.~1256--1300, the exact dating is debated). The Latin version
became known as \textit{Picatrix}, by which name it has been
transmitted ever since.

\noindent\textbf{The surviving Latin manuscripts.} \citet{pingree1986},
in the manuscript census of the introduction to his critical edition,
documents approximately twenty surviving Latin Picatrix manuscripts
of varying completeness, dated between the late thirteenth and the
sixteenth century. The Cracow Jagiellonian Library MS~793 is the only
one of these that preserves a fully developed illustrative program,
which is why it occupies a privileged place in the iconographic
analysis of \citet{sniezynska2009} and in the present paper.

\noindent\textbf{Modern editions.} The principal modern critical
edition of the Latin text is \citet{pingree1986}; the principal
French translation with commentary is \citet{bakhouche2003}; the
principal English translation is \citet{attrell2019}, working from
Pingree's Latin and incorporating insights from the Arabic
\textit{Gh\=ayat al-\d{H}ak\={\i}m}. All three are cited in the
present paper as needed.

\subsection{The \textit{karacteres} prescription}
\label{sec:karacteres_prescription}

Within the Picatrix text, the specific category that matters for the
Voynich comparison is the \textit{karacteres} prescription. The
Picatrix instructs its operator, in many of its planetary-image
recipes, to inscribe specific images (statues, plates, rings, combs,
amulets) with reserved non-alphabetic characters drawn from a
restricted graphic vocabulary. \citet{sniezynska2009}~(p.~103)
describes these characters succinctly:

\begin{quote}
\itshape
``[The karacteres are] transformations of the Arabic or Semitic
alphabet, encountered as medieval mason's marks.''
\end{quote}

\noindent The defining properties of the category, as
\citet{sniezynska2009} reconstructs them, are:

\begin{itemize}[itemsep=0pt,topsep=2pt]
\item the inventory is small and closed (a finite set of geometric
forms, not an open writing system);
\item the forms are constructed from a restricted graphic vocabulary
of straight strokes, small circles, right-angle hooks, and curved
tails (the same vocabulary documented in
\S\ref{sec:vc2_shields}--\S\ref{sec:vc4_alphabet});
\item the function is operative, not communicative: the inscription
is a component of a talismanic device whose efficacy depends on the
correct deployment of the characters in the prescribed iconographic
context;
\item the characters are explicitly prescribed for specific
planetary positions, with the Sun and Venus chapters of Book~II
preserving the most extensive surviving instances (cf.\
\S\S\ref{sec:vc2_shields}--\ref{sec:oloiol}).
\end{itemize}

\noindent The same author also documents (p.~105) the historical fate
of the karacteres tradition in early-modern scholastic reception:

\begin{quote}
\itshape
``The use of the magical alphabet and esoteric symbols is condemned
in the \textit{Speculum astronomiae}; perhaps this is why such
inscriptions contained in the text of \textit{Picatrix Latinus}
appear in the miniatures of the Cracow manuscript only once.''
\end{quote}

\noindent The reference to the \textit{Speculum astronomiae}
(thirteenth-century, traditionally attributed to Albertus Magnus)
captures the doctrinal pressure that led the Cracow scribe of
1471--1474 to suppress most of the karacteres prescribed by his Latin
source text, preserving only three vestiges (the two Sun-shield
sigils of f.~191r and the $\mathrm{O\Lambda OIO\Lambda}$ Venus
inscription of f.~191v) on iconographic grounds. The asymmetry of
preservation that this scribal censorship produces is exactly the
asymmetry that \S\ref{sec:more_complete} exploits to argue that the
Voynich, on the karacteres dimension specifically, is a more
complete witness of the medieval tradition than its own documented
Cracow contemporary.

The \textit{karacteres} prescription is therefore the SPECIFIC textual
category---closed inventory, restricted graphic vocabulary, operative
function, planetary association, doctrinally censored in the
late-medieval scholastic reception---into which the empirical findings
of \S\S\ref{sec:subalphabet}--\ref{sec:two_phenomena} most naturally
fit. The interpretive proposal of \S\ref{sec:reformulation}, that
Voynich f57v's Phenomenon~B is one more witness of the same
documented tradition, is the obvious reading of those findings under
the theoretical frame established here. The remaining sections develop
that reading and document its limitations.


% ============================================================
\section{Quantitative finding: the private subalphabet of f57v}
\label{sec:subalphabet}

This section is the empirical core of the paper. Working only from the
public STA1~2.0 transliteration of the Voynich Manuscript and from the
locus-tagging conventions of the IVTFF format \citep{zandbergen_sta1,zandbergen_ivtff},
we show that folio~f57v carries a textually anomalous inventory: five glyphs
that occur \emph{nowhere else} in the entire 224-folio corpus, plus three
further glyphs that are concentrated on f57v at enrichment factors of two to
three orders of magnitude over the corpus baseline. The probability of this
concentration arising by chance is on the order of~$10^{-26}$ for the single
most frequent exclusive glyph and~$10^{-64}$ for the five exclusive glyphs
considered jointly. Within the folio, the anomaly is further localized to the
four concentric text rings of the diagram---most densely to ring~f57v.3---and
is entirely absent from the central nymph labels. The full computation is
released as the public companion notebook~74
(\texttt{74\_f57v\_private\_subalphabet.ipynb}) and as the structured JSON
results file \texttt{f57v\_subalphabet\_resultados.json}.

\subsection{Methodology and metrics}
\label{sec:method}

We work from the STA1~2.0 super-transliteration of the Voynich Manuscript
\citep{zandbergen_sta1}: 37,087 word tokens, 157,254 atom occurrences, 166
distinct atom types distributed over 224 folios. STA1 uses a finer-grained
inventory than the more widely cited EVA transliteration; the additional
distinctions matter for this analysis, because several of the glyphs we will
identify as exclusive to f57v are EVA-collapsed into more common atoms in EVA
(see~\S\ref{sec:eva_caveat}).

For every atom type~$a$ in the STA1 inventory we compute its global frequency
$N(a)$ across the 224-folio corpus and its frequency $n_{f}(a)$ on each
individual folio~$f$. The \emph{folio share} of an atom on folio~$f$ is
$n_{f}(a)/N(a)$, and its \emph{enrichment factor} on folio~$f$ is the folio
share divided by the folio's atom share of the corpus. For f57v, which
contains $428$~atoms out of $154{,}973$~corpus atoms (after removing
high-uncertainty IVTFF tokens), the atom share is $0.276\%$, so the baseline
expectation for any atom under a uniform distribution is that $0.276\%$ of its
total occurrences should land on f57v. Atoms whose actual concentration on
f57v is dramatically above this baseline are the candidates of interest.

A second metric is the per-line proportion of \emph{isolated atoms}:
single-character word tokens that have no neighbours within the same word.
A locus written as a list of one-character ``words'' separated by spaces is
the IVTFF representation of an enumerative inscription (an alphabet table, a
list of separate sigils, a label key) as opposed to running text. The
proportion of isolated atoms per folio is therefore a coarse but transcription-stable
proxy for ``how much of this folio looks like a list rather than a paragraph''.

Finally, we use the IVTFF locus tag system \citep{zandbergen_ivtff} to identify
which lines of f57v belong to which page element: \texttt{@L0} marks centred
labels (the central nymph captions of f57v), \texttt{@Ro} marks rotated radial
labels, and \texttt{@Cc}/\texttt{+Cc} mark concentric circular text. This
distinction is essential because the anomaly is page-element-specific, not
folio-uniform.

\subsection{Identification of exclusive and enriched glyphs}
\label{sec:identification}

Direct enumeration of per-glyph folio shares yields a small but extremely
sharp population of anomalies. Five STA1 atoms have $100\%$ of their corpus
occurrences on f57v, and three further atoms have folio shares between
$20\%$ and $84\%$---all hundreds of times the $0.276\%$ baseline.
Table~\ref{tab:subalphabet} lists the eight glyphs together with their global
counts, f57v counts, percentage on f57v, and enrichment factor over baseline.

\begin{table}[h]
\centering
\small
\begin{tabular}{lrrrr}
\toprule
Glyph & Global $N(a)$ & On f57v $n_{f57v}(a)$ & \% on f57v & Enrichment \\
\midrule
\multicolumn{5}{l}{\emph{Exclusive (folio share = 100\%):}} \\
X2 & 10 & 10 & 100.0\% & $362\times$ \\
Xd & 5  & 5  & 100.0\% & $362\times$ \\
Pc & 5  & 5  & 100.0\% & $362\times$ \\
Xf & 4  & 4  & 100.0\% & $362\times$ \\
Ea & 1  & 1  & 100.0\% & $362\times$ \\
\midrule
\multicolumn{5}{l}{\emph{Strongly enriched:}} \\
Eb & 6  & 5  & 83.3\%  & $302\times$ \\
Ja & 14 & 4  & 28.6\%  & $103\times$ \\
X1 & 34 & 7  & 20.6\%  & $\phantom{0}75\times$ \\
\bottomrule
\end{tabular}
\caption{The eight atoms forming the f57v private subalphabet. Counts are
measured against the STA1~2.0 transliteration of the entire 224-folio Voynich
corpus. f57v's atom share of the corpus is $0.276\%$; an atom occurring at
the corpus baseline rate would therefore have a folio share of $0.276\%$, and
the enrichment factor is the folio share divided by this baseline. Source:
\texttt{74\_f57v\_private\_subalphabet.ipynb}.}
\label{tab:subalphabet}
\end{table}

The five exclusive glyphs collectively contribute $25$ atom tokens, and the
three enriched glyphs contribute a further $16$ tokens on f57v. These $41$
tokens make up nearly $10\%$ of the $428$ atoms on the folio---a substantial
fraction of the page is written in glyphs that essentially do not appear
elsewhere in the manuscript.

\subsection{Statistical significance and probability of chance}
\label{sec:probability}

A simple null model treats the corpus as $428/154{,}979 = 0.276\%$ ``chance
of any given atom landing on f57v''. Under this null, the probability that
all $10$~tokens of glyph X2 land on f57v by chance alone is
\[
p_{X2} \;=\; (0.00276)^{10} \;\approx\; 2.6\times 10^{-26}.
\]
Computing the same probability for each of the five exclusive glyphs and
multiplying gives a combined null probability for the $25$ exclusive tokens
of approximately $1.1\times 10^{-64}$. This is the probability under a
strict ``random sprinkling'' null and is therefore an \emph{upper bound} on
the chance probability under any reasonable urn model that does not already
encode some form of folio-level concentration. Even granting several orders
of magnitude of slack for non-uniformities in the corpus (different folios
have different sizes, some glyphs cluster topically, etc.), the f57v
concentration is statistically incompatible with chance by tens of orders of
magnitude.

The interpretation is unambiguous: the eight glyphs of
Table~\ref{tab:subalphabet} are not a random fluctuation. They are a textually
distinctive sub-inventory used by the scribe \emph{exclusively or near-exclusively}
on f57v, and we will refer to them collectively as the \emph{f57v private
subalphabet}.

\subsection{Intra-folio distribution: where the subalphabet lives}
\label{sec:intrafolio}

Folio f57v is itself structurally heterogeneous: it contains $13$ text loci
(four central nymph labels, four radial nymph labels, four concentric text
rings, and one further short label). Table~\ref{tab:perline} breaks down the
$428$ atoms of the folio by IVTFF locus tag, giving for each line its total
atom count, word count, isolated-atom count, and the breakdown of subalphabet
occurrences.

\begin{table}[h]
\centering
\small
\begin{tabular}{llrrrl}
\toprule
Locus & Tag & Atoms & Words & Isolated & Subalphabet glyphs \\
\midrule
f57v.1  & \texttt{@L0} & 5   & 1  & 0  & --- \\
f57v.2  & \texttt{@Cc} & 148 & 50 & 17 & X2$\times$3 \\
f57v.3  & \texttt{+Cc} & 68  & 68 & 68 & X2$\times$4, X1$\times$4, Xf$\times$4, Xd$\times$4, Pc$\times$4, Ea$\times$1, Eb$\times$4, Ja$\times$3 \\
f57v.4  & \texttt{+Cc} & 89  & 31 & 11 & X2$\times$2, X1$\times$1, Eb$\times$1 \\
f57v.5  & \texttt{+Cc} & 57  & 35 & 28 & X2$\times$1, X1$\times$2, Xd$\times$1, Pc$\times$1, Ja$\times$1 \\
f57v.6  & \texttt{@L0} & 8   & 1  & 0  & --- \\
f57v.7  & \texttt{@L0} & 9   & 1  & 0  & --- \\
f57v.8  & \texttt{@L0} & 7   & 1  & 0  & --- \\
f57v.9  & \texttt{@L0} & 8   & 1  & 0  & --- \\
f57v.10 & \texttt{@Ro} & 7   & 1  & 0  & --- \\
f57v.11 & \texttt{@Ro} & 8   & 1  & 0  & --- \\
f57v.12 & \texttt{@Ro} & 5   & 1  & 0  & --- \\
f57v.13 & \texttt{@Ro} & 9   & 2  & 0  & --- \\
\bottomrule
\end{tabular}
\caption{Per-line breakdown of folio f57v by IVTFF locus tag. The four
concentric text rings (\texttt{@Cc}/\texttt{+Cc}, lines f57v.2--f57v.5) carry
the entire subalphabet. The eight central labels (\texttt{@L0}/\texttt{@Ro})
contain zero subalphabet atoms. Line~f57v.3 alone accounts for $28$ of the
$41$ subalphabet tokens on the folio and is the only line in the entire
corpus where every word is a single-character isolated atom ($68/68 = 100\%$).
Source: \texttt{74\_f57v\_private\_subalphabet.ipynb}.}
\label{tab:perline}
\end{table}

Three properties of the table are worth emphasising. First, the subalphabet
is \emph{entirely confined to the four concentric text rings}: the central
nymph labels (\texttt{@L0}) and the radial labels (\texttt{@Ro}) of the same
folio contain none of the eight glyphs. The anomaly is therefore not a
property of f57v as such, but of the ring text within f57v.

Second, line~f57v.3 is qualitatively different from every other line in the
manuscript. All $68$ of its ``words'' are single-atom tokens; this is the
only line in the entire 224-folio corpus where the isolated-atom rate is
exactly $100\%$. Of the eight subalphabet glyphs, seven appear here---five of
them in groups of exactly four occurrences. The structure of f57v.3 is that
of an enumeration, not a sentence: $17$ atomic characters repeated $4$ times.
This $17\!\times\!4$ regularity has been noted intuitively by voynichologists
for decades; what is new here is the quantitative demonstration that the
characters in question are themselves a closed inventory used essentially
nowhere else in the manuscript.

Third, all four ring lines (f57v.2--f57v.5) carry subalphabet atoms in
descending density: 12 on f57v.2, 28 on f57v.3, 4 on f57v.4, 6 on f57v.5.
The subalphabet is therefore a property of the ring system as a whole, with
f57v.3 as its densest concentration.

\subsection{Comparison with the 224-folio corpus and rank ordering}
\label{sec:ranking}

To put f57v's anomaly in context against the rest of the manuscript, we rank
all $224$ folios by their proportion of isolated atoms. Table~\ref{tab:ranking}
shows the top~15. The median across the corpus is $2.3\%$.

\begin{table}[h]
\centering
\small
\begin{tabular}{rlrrrr}
\toprule
Rank & Folio & Words & Atoms & Avg. word len. & \% isolated atoms \\
\midrule
1  & \textbf{f57v}  & 194 & 428  & 2.21 & \textbf{63.92\%} \\
2  & f49v   & 154 & 589  & 3.82 & 20.13\% \\
3  & f38v   & 58  & 205  & 3.53 & 12.07\% \\
4  & f66r   & 338 & 1406 & 4.16 & 10.65\% \\
5  & f25v   & 57  & 208  & 3.65 & 8.77\% \\
6  & f90r2  & 46  & 171  & 3.72 & 8.70\% \\
7  & f93r   & 151 & 600  & 3.97 & 8.61\% \\
8  & f35r   & 84  & 313  & 3.73 & 8.33\% \\
9  & f90v2  & 62  & 265  & 4.27 & 8.06\% \\
10 & f38r   & 38  & 147  & 3.87 & 7.89\% \\
11 & f6v    & 107 & 399  & 3.73 & 7.48\% \\
12 & f24r   & 108 & 361  & 3.34 & 7.41\% \\
13 & f21r   & 97  & 377  & 3.89 & 7.22\% \\
14 & f42v   & 97  & 392  & 4.04 & 7.22\% \\
15 & f67r1  & 163 & 656  & 4.02 & 6.75\% \\
\bottomrule
\end{tabular}
\caption{Top~15 folios by percentage of isolated atoms. f57v is rank~\#1 at
$63.9\%$, more than three times rank~\#2 and roughly $28\times$ the corpus
median ($2.3\%$). Ranks~\#2 and~\#4 (f49v and f66r) are documented in
\S7.2 of the IVTFF specification \citep{zandbergen_ivtff} as folios with
extensive ``key-like'' transcription artefacts (single-character entries
along page margins) that partially inflate their isolated-atom counts.
Source: \texttt{74\_f57v\_private\_subalphabet.ipynb}.}
\label{tab:ranking}
\end{table}

f57v is rank~\#1 by a wide margin: its isolated-atom rate of $63.9\%$ is
more than three times that of the second-place folio and approximately
$28$~times the corpus median. Ranks~\#2 and~\#4 (f49v at $20.1\%$ and f66r
at $10.6\%$) are flagged in the IVTFF specification \citep{zandbergen_ivtff}
as folios with marginal ``key-like'' transcription artefacts: f49v has a
sequence of single-character entries along the right margin that the
transcriber rendered as isolated atoms, and f66r contains a multi-line
list-like marginal annotation. After mentally discounting those two folios
on transcription grounds, the next genuine running-text folio is f38v at
$12.1\%$. Even so, f57v at $63.9\%$ is approximately $5\times$ further from
its nearest comparator. \textbf{f57v is the only folio in the manuscript
whose dominant text behaviour is enumerative rather than continuous}, and it
is so by a margin that places it in a class of its own.

\subsection{A note on transliteration: STA1 versus EVA}
\label{sec:eva_caveat}

The five exclusive atoms reported here (X2, Xd, Xf, Pc, Ea) and the three
enriched atoms (Eb, Ja, X1) are STA1 designations. Several of these
correspond in EVA to glyphs that are merged with more common atoms: for
example, several of the ``X-series'' STA1 atoms collapse into the EVA atoms
\texttt{x} or \texttt{r}. Repeating the analysis in EVA does not eliminate
the anomaly---f57v is still rank~\#1 by isolated-atom percentage in EVA---but
it does \emph{flatten} the table of exclusive atoms: where STA1 separates a
genuinely f57v-exclusive shape from a more common look-alike, EVA collapses
them, and the apparent exclusivity disappears under that conflation. The
quantitative result of this section therefore depends on the finer
inventory of STA1, not on EVA. We regard this as evidence in favour of
STA1's distinctions being substantive rather than arbitrary, at least on
this folio. The same conclusion was reached for unrelated reasons by the
transcription robustness analysis of the author's first companion paper
\citep{jimenez2026negative}.


% ============================================================
\section{Robustness: the cosmological-suite test}
\label{sec:f70r}

The single most important threat to the interpretation of \S\ref{sec:subalphabet}
is the possibility that the f57v subalphabet is not a property of f57v as such
but of an entire \emph{class} of cosmological pages within the manuscript. If
that were the case, the f57v anomaly would lose its specificity: the
``private'' subalphabet would in reality be the dialect of a wider section of
the codex, and the link with a single talismanic inscription on Venus's comb
would be much weaker. To rule this out we test the subalphabet against every
other folio in the cosmological/astronomical block of the manuscript, not only
against the most obvious iconographic sibling.

\subsection{The cosmological block of the Voynich and the test set}
\label{sec:cosmoblock}

The cosmological/astronomical block of the Voynich consists of f57v itself
together with the foldout pages of quires~9 and~10: f67r1, f67r2, f67v1,
f67v2, f68r1, f68r2, f68r3, f68v1, f68v2, f68v3, and the f70r and f70v solar
foldout pages (f70r1, f70r2, f70v1, f70v2). All fourteen of these pages share
the broad iconographic family of f57v: a circular or radial diagram, often
with a central anthropomorphic figure (Sun, Moon, or human face), surrounded
by concentric rings of stars, glyphs and text labels. Several of them
(f67r1, f68r1, f68v3, f70r1) display an explicit anthropomorphic sun or moon
face at the centre, exactly as f57v's $\odot$ symbol prefigures schematically.

The natural null hypothesis is therefore: \emph{if the eight subalphabet
glyphs of Table~\ref{tab:subalphabet} are characteristic of the cosmological
genre rather than of f57v specifically, they should appear with appreciable
frequency on at least some of these fourteen other folios}. This is a much
stricter test than comparing f57v to a single ``twin'' folio, because the
test set spans the entire cosmological iconographic family of the manuscript.

\subsection{Test results}
\label{sec:test_results}

We computed the per-folio occurrence count of each of the eight subalphabet
glyphs across the fourteen test folios using the same STA1 corpus and the
same enumeration procedure as \S\ref{sec:subalphabet}. The results are shown
in Table~\ref{tab:cosmotest}.

\begin{table}[h]
\centering
\small
\begin{tabular}{lrrrrrrrrrr}
\toprule
Folio & Atoms & X1 & X2 & Xd & Xf & Pc & Ea & Eb & Ja & Total \\
\midrule
\textbf{f57v}  & 428 & 7 & 10 & 5 & 4 & 5 & 1 & 5 & 4 & \textbf{41} \\
\midrule
f67r1 & 656 & 0 & 0 & 0 & 0 & 0 & 0 & 0 & 0 & 0 \\
f67r2 & 749 & 0 & 0 & 0 & 0 & 0 & 0 & 0 & 0 & 0 \\
f67v1 & 336 & 0 & 0 & 0 & 0 & 0 & 0 & 0 & 1 & 1 \\
f67v2 & 276 & 0 & 0 & 0 & 0 & 0 & 0 & 0 & 0 & 0 \\
f68r1 & 317 & 0 & 0 & 0 & 0 & 0 & 0 & 0 & 1 & 1 \\
f68r2 & 391 & 0 & 0 & 0 & 0 & 0 & 0 & 0 & 0 & 0 \\
f68r3 & 476 & 0 & 0 & 0 & 0 & 0 & 0 & 0 & 0 & 0 \\
f68v1 & 381 & 0 & 0 & 0 & 0 & 0 & 0 & 0 & 0 & 0 \\
f68v2 & 455 & 0 & 0 & 0 & 0 & 0 & 0 & 0 & 0 & 0 \\
f68v3 & 692 & 0 & 0 & 0 & 0 & 0 & 0 & 0 & 0 & 0 \\
f70r1 & 495 & 0 & 0 & 0 & 0 & 0 & 0 & 0 & 0 & 0 \\
f70r2 & 925 & 0 & 0 & 0 & 0 & 0 & 0 & 0 & 0 & 0 \\
f70v1 & 400 & 0 & 0 & 0 & 0 & 0 & 0 & 0 & 0 & 0 \\
f70v2 & 585 & 0 & 0 & 0 & 0 & 0 & 0 & 0 & 0 & 0 \\
\midrule
\textbf{Sum (14 cosmo)} & 7134 & \textbf{0} & \textbf{0} & \textbf{0} & \textbf{0} & \textbf{0} & \textbf{0} & \textbf{0} & \textbf{2} & \textbf{2} \\
\bottomrule
\end{tabular}
\caption{Subalphabet occurrence counts across f57v and all 14 other
cosmological/astronomical folios of the Voynich Manuscript. Out of $7134$
atoms across the 14 test folios, the five exclusive glyphs (X2, Xd, Xf, Pc,
Ea) and the two most-enriched non-exclusive glyphs (X1, Eb) all score
\emph{exactly zero}. Only the weakest enriched glyph, Ja, contributes two
isolated stray tokens (one each on f67v1 and f68r1). Source:
\texttt{f57v\_subalphabet\_cosmological\_robustness.json}.}
\label{tab:cosmotest}
\end{table}

The result is unambiguous. Of the eight subalphabet glyphs, seven do not
appear at all on any of the fourteen other cosmological folios:
\textbf{0 occurrences across $7134$~atoms of comparison material}. The only
exception is glyph Ja, which contributes one token on f67v1 and one on
f68r1---but Ja is also the weakest member of the subalphabet (folio share
$28.6\%$ on f57v, the lowest of the eight) and the two stray tokens are
fully consistent with its global distribution outside f57v ($14$~tokens
total in the corpus, of which $4$ on f57v and $10$ scattered elsewhere).

\subsection{What the cosmological-suite test rules out}
\label{sec:f70r_meaning}

Three alternative explanations of the f57v anomaly are decisively ruled out
by Table~\ref{tab:cosmotest}:

\noindent (i)~\textbf{Genre artefact.} ``f57v's vocabulary is the vocabulary
of the cosmological-page genre.'' If this were true, the same vocabulary
should appear on at least some of f67r1, f68r1, f70r1 (the other diagrams
with explicit central solar figures). It appears on \emph{none} of them.

\noindent (ii)~\textbf{Foldout artefact.} ``Foldout pages have larger writing
surfaces and admit more glyph variation than ordinary pages.'' If this were
true, the foldouts of ff.~67--68 and 70 should sample at least the most
common subalphabet glyphs (X1 with 34 corpus occurrences, Ja with 14, X2 with
10). None of them does, and the foldouts collectively contain $7134$~atoms
of comparison material---more than $16\times$ f57v's atom count---across
which the subalphabet has zero footprint.

\noindent (iii)~\textbf{Diagram artefact.} ``Any folio with concentric text
rings will produce isolated single-character tokens.'' If this were true,
f70r1 (an explicit anthropomorphic sun diagram with 16 radial rays and
labels) and the f67/f68 diagrams (concentric star rings) should show some
elevation of the isolated-atom rate. None do (cf.\ Table~\ref{tab:ranking};
f70r1 sits at $1\%$ isolated atoms, f70r2 at $3\%$, all f67/f68 foldouts
likewise within ordinary corpus range).

The conclusion is that the f57v subalphabet is \textbf{folio-specific, not
genre-specific, not foldout-specific, and not diagram-specific}. It is a
property of the four concentric text rings of f57v in particular, and within
those rings most densely a property of line f57v.3. This is the strongest
possible robustness test that could have failed---the test set is the entire
cosmological iconographic family of the manuscript---and it does not fail.


% ============================================================
\section{The unified categorial labeling system of the Voynich}
\label{sec:unified}

The previous two sections have isolated f57v as textually unique within the
manuscript: only f57v carries the eight-glyph private subalphabet
(\S\ref{sec:subalphabet}), and only f57v has its concentric rings filled
with isolated single-character ``words'' (\S\ref{sec:f70r}). One might infer
from this uniqueness that f57v is altogether disconnected from the rest of
the manuscript---an alien folio bound into the codex by accident, written
in a separate hand or system. The purpose of this section is to show that
this inference would be wrong. The \emph{labels} of f57v---the central
nymph captions and the radial captions, written in normal Voynich glyphs
rather than in subalphabet glyphs---turn out to follow the same canonical
positional morphology that the author's previous unpublished work
documents for the zodiac block of the manuscript (notebooks~67--71). The
same is true of the radial labels on f70r1 (the sun foldout). Voynich f57v
is therefore not a textual outlier but a folio with two coexisting
textual layers: a normal labeling layer (analysed in this section) and a
localized inscription layer (the subject of \S\ref{sec:subalphabet} and
\S\ref{sec:f70r}). The two layers will be named explicitly as
Phenomenon~A and Phenomenon~B in \S\ref{sec:two_phenomena}.

\subsection{The zodiac canonical morphology (background)}
\label{sec:zodiac_canon}

The zodiac block of the Voynich Manuscript (folios f70v2--f73v) contains
ten signs of the zodiac drawn as roundels surrounded by 30~degree-marker
nymphs each, every nymph carrying a small radial label of typically 4--8
glyphs. The author's previous unpublished work on this block (notebooks
67--71 of the project repository) established the following structural
properties of these labels:

\begin{enumerate}[itemsep=0pt,topsep=2pt]
\item There are 27~canonical zodiac labels recurring across the block.
\item All 27~canonical labels follow the morphology
\[
[\,A1\,] \;+\; [\,Q1/Q2\,] \;+\; [\,J\text{-class season encoder}\,] \;+\;
[\,\text{categorial body}\,] \;+\; [\,A2\,]
\]
\item In position~1, $23/27 = 85\%$ of labels begin with the atom~A1; the
remaining four positions are taken by structurally similar atoms (A3, K1,
Aa).
\item In position~2, the slot is dominated by the binary
$\{\mathrm{Q1}, \mathrm{Q2}\}$, which the analysis of notebook~71
identifies as a season encoder: this slot is significantly correlated
with the calendrical position of the nymph within the sign
($p = 0.0002$, Bonferroni-corrected over the four positional tests).
\item Positions~3 and~4 carry categorial atoms that we interpret as
``fertility'' and ``sign type'' encoders respectively (notebook~71
section~5).
\item A final A2 marker closes the canonical word in $\sim 80\%$ of
labels.
\end{enumerate}

The full statistical defence of these properties is the subject of a
separate paper now in preparation, and we cite it here only as
background. For the present argument, the only thing we need is the
observation that there is a \emph{canonical positional template} for
zodiac labels and that the template is empirically tight enough to be
falsifiable on new label sets.

\subsection{Test~1: positional analysis on f57v central nymph labels}
\label{sec:test_l0}

Folio f57v contains four central nymph labels, one per nymph, tagged
\texttt{@L0} in the IVTFF transliteration. After tokenisation into STA1
atoms, the four labels decompose as shown in Table~\ref{tab:l0}.

\begin{table}[h]
\centering
\small
\begin{tabular}{lllllllllll}
\toprule
Locus & p1 & p2 & p3 & p4 & p5 & p6 & p7 & p8 & p9 & length \\
\midrule
f57v.6 & A1 & Q2 & A1 & B1 & A3 & C1 & A3 & B4 &     & 8 \\
f57v.7 & A1 & P1 & A3 & C1 & A3 & F3 & B1 & B2 & A2 & 9 \\
f57v.8 & A1 & B2 & Q1 & K1 & B1 & A3 & B2 &    &     & 7 \\
f57v.9 & A1 & Q2 & A3 & C1 & B1 & A3 & B2 & A2 &     & 8 \\
\bottomrule
\end{tabular}
\caption{Positional decomposition of the four central nymph labels of
f57v (\texttt{@L0}) into STA1 atoms. All four begin with A1 in
position~1 ($4/4 = 100\%$, vs.\ $85\%$ in the zodiac canon); three of
four contain a Q-class atom (Q1 or Q2) in positions~2--3; A3 appears in
all four labels in either position~3 or position~5; and the closing A2
marker is present in $2/4$ labels (f57v.7, f57v.9). Source:
\texttt{74\_f57v\_private\_subalphabet.ipynb} key
\texttt{nymph\_L0\_positional}.}
\label{tab:l0}
\end{table}

The fit with the zodiac canonical morphology is exact in position~1 and
strong in positions~2--3:

\noindent\textbf{Position 1.} All four labels start with A1. The zodiac
distribution at position~1 over the 27~canonical labels is
$\{A1\!:\!23, A3\!:\!1, K1\!:\!2, Aa\!:\!1\}$, i.e.\ $85\%$ A1. The
nymph labels of f57v therefore conform to the dominant zodiac value with
$100\%$ purity, more strictly than the zodiac block itself. Under a
naive null hypothesis where each label's position-1 atom is sampled
uniformly from the corpus inventory of $\sim 166$~atom types, the
probability of $4/4$ labels independently starting with the same
specific atom is $(1/166)^3 \approx 2.2\times 10^{-7}$.

\noindent\textbf{Positions 2--3.} Of the four labels, three contain a
Q-class atom (Q1 or Q2) in either position~2 or position~3 (f57v.6,
f57v.8, f57v.9), as the zodiac canon prescribes for the season encoder
slot. The fourth label (f57v.7) places P1 in position~2 instead, but P1
is structurally similar to Q1/Q2 (it shares the same allograph cluster
in the unpublished allograph analysis of nb27). The zodiac season-slot
discipline is therefore satisfied in 4/4 labels, three strictly and one
loosely.

\noindent\textbf{Categorial body.} All four labels contain at least one
A3 atom in positions~3--5 and at least one B-class atom (B1, B2, or
B4) in positions~5--7. This matches the categorial-body slots of the
zodiac canon.

\noindent\textbf{Closing A2 marker.} Two of the four labels (f57v.7 and
f57v.9) end with A2; the other two terminate with a B-class atom. This
$50\%$ A2 closure is below the zodiac canon's $\sim 80\%$ but is
consistent with it given the small sample.

The four central nymph labels of f57v therefore conform to the
zodiac canonical morphology in essentially all of its discriminative
slots. They are not random Voynich words; they are structurally zodiac
labels, even though they appear on a non-zodiac folio.

\subsection{Test~2: positional analysis on f70r1 sun-ray labels}
\label{sec:test_ri}

Folio f70r1 (the sun foldout, with an anthropomorphic sun face at the
centre and 16 radial rays) carries 15 short radial labels tagged
\texttt{@Ri} in the IVTFF transliteration, one per ray. We applied the
same canonical-template test to these 15 labels. The aggregate result is
shown in Table~\ref{tab:ri}.

\begin{table}[h]
\centering
\small
\begin{tabular}{lr}
\toprule
Property & \% of labels (n=15) \\
\midrule
Begin with canonical position-1 atom (A1 or A2)        & 86.7\% (13/15) \\
Contain Q1 or Q2 in early positions                    & 73.3\% (11/15) \\
Contain J-class season encoder (J1)                    & 33.3\% (5/15) \\
Contain any of the eight subalphabet glyphs of \S\ref{sec:subalphabet} & \textbf{0.0\% (0/15)} \\
\bottomrule
\end{tabular}
\caption{Aggregate canonical-morphology compliance of the 15 sun-ray
labels (\texttt{@Ri}) on folio f70r1. The conformance with the zodiac
canon is essentially the same as on f57v's central nymph labels, while
the subalphabet of \S\ref{sec:subalphabet} is entirely absent. Source:
\texttt{74\_f57v\_private\_subalphabet.ipynb} key \texttt{f70r1\_Ri}.}
\label{tab:ri}
\end{table}

The compliance rates with positions~1 and~2 of the zodiac canon
($87\%$ and $73\%$) are well within the zodiac block's own observed
ranges ($85\%$ and $\sim 70\%$ respectively). The J-class season
encoder appears at $33\%$, which is lower than the zodiac block's
$\sim 60\%$ but still substantially above the corpus-baseline rate
for J1.

The most striking observation, however, is the last row: among the
15~radial labels of f70r1, exactly \textbf{zero} contain any of the
eight subalphabet glyphs of \S\ref{sec:subalphabet}. The labels of
f70r1 are written in fully canonical Voynich vocabulary, despite
sitting on a folio that is iconographically the closest sibling of
f57v. This is consistent with the cosmological-suite test of
\S\ref{sec:f70r} and reinforces the same conclusion from a different
angle: the canonical labeling system extends across cosmological
folios, while the f57v subalphabet does not.

\subsection{The cross-folio shared word A1Q1J1A3C1}
\label{sec:shared_word}

A particularly clear instance of cross-folio canonical morphology is
a recurring word, which we will refer to as \textbf{A1\,Q1\,J1\,A3\,C1}.
Decomposed positionally:

\begin{itemize}[itemsep=0pt,topsep=2pt]
\item position~1: A1 (canonical zodiac p1)
\item position~2: Q1 (canonical zodiac p2 season encoder)
\item position~3: J1 (canonical zodiac p3 J-class encoder)
\item position~4: A3 (canonical zodiac p4 categorial atom)
\item position~5: C1 (canonical zodiac p5 closing slot)
\end{itemize}

This is, slot for slot, a complete canonical zodiac label of the form
described in \S\ref{sec:zodiac_canon}. Across the entire STA1 corpus
the word A1Q1J1A3C1 occurs ten times, in exactly the loci listed in
Table~\ref{tab:shared_word}.

\begin{table}[h]
\centering
\small
\begin{tabular}{llll}
\toprule
Locus & Folio & Tag & Folio context \\
\midrule
f31v.9    & f31v   & \texttt{+P0} & herbal section \\
f57v.13   & f57v   & \texttt{@Ro} & central radial label of f57v \\
f70r2.15  & f70r2  & \texttt{@Cc} & solar foldout (concentric ring) \\
f72v3.20  & f72v3  & \texttt{@Cc} & zodiac block (Leo) \\
f88v.9    & f88v   & \texttt{+P0} & pharmaceutical section \\
f102v2.21 & f102v2 & \texttt{+P0} & pharmaceutical section \\
f103v.9   & f103v  & \texttt{+P0} & recipe section \\
f106v.16  & f106v  & \texttt{+P0} & recipe section \\
f108v.15  & f108v  & \texttt{+P0} & recipe section \\
f111v.9   & f111v  & \texttt{+P0} & recipe section \\
\bottomrule
\end{tabular}
\caption{All ten occurrences of the canonical-zodiac-form word
\textbf{A1Q1J1A3C1} in the STA1 corpus. The word appears once on f57v
itself (as a central radial label), once on the f70r2 solar foldout (as
ring text), once on the f72v3 zodiac folio (Leo), and seven times in the
herbal/pharmaceutical/recipe sections. Source:
\texttt{74\_f57v\_private\_subalphabet.ipynb} key
\texttt{shared\_word\_A1Q1J1A3C1}.}
\label{tab:shared_word}
\end{table}

Three observations from Table~\ref{tab:shared_word}:

\noindent (i)~The word appears in three of the four major sections of
the manuscript (herbal, cosmological/zodiac, pharmaceutical, recipes).
It is not section-specific. This is the strongest single piece of
evidence that the canonical zodiac morphology is in fact a
\emph{manuscript-wide} positional system, not a property of the
zodiac block alone.

\noindent (ii)~The cosmological occurrences (f57v.13, f70r2.15) and
the zodiac occurrence (f72v3.20) place the same canonical word in the
same canonical syntactic context---a label or short ring entry,
typically associated with a circular diagram. The categorial system
treats f57v's radial labels and the Leo nymphs of f72v3 as
\emph{the same kind of textual element}.

\noindent (iii)~The seven occurrences of A1Q1J1A3C1 in the herbal,
pharmaceutical and recipe sections (where they appear in
\texttt{+P0} continuous-text loci rather than as labels) suggest that
the canonical zodiac morphology is also embedded in the running text
of those sections. This is consistent with the encyclopaedic-code
hypothesis of the second companion paper~\citep{jimenez2026twotier}:
the canonical morphology is the manuscript's general categorial
labeling system, and labels-as-such are simply the most visible
instances of it.

The combined picture is clear. The canonical zodiac positional template
is not a local feature of the zodiac block; it operates on the central
nymph labels of f57v, on the radial labels of f70r1, on the ring text
of f70r2, and even in the body text of the herbal and pharmaceutical
sections. f57v participates fully in this manuscript-wide system---and
this is the textual layer of f57v that we will call Phenomenon~A in
the next section.


% ============================================================
\section{Two textually distinct phenomena coexist in the Voynich}
\label{sec:two_phenomena}

The combined picture of \S\S\ref{sec:subalphabet}--\ref{sec:unified} forces a
conceptual reframing of the manuscript. The Voynich contains \emph{two}
textually distinct phenomena that coexist in the same codex---and, in the
case of f57v, on the same folio---without contaminating each other. Calling
them ``Phenomenon~A'' and ``Phenomenon~B'' is a deliberate choice of neutral
terminology, intended to keep the empirical separation distinct from any
particular interpretive hypothesis. The interpretive proposal (Phenomenon~A
as descriptive categorial labeling, Phenomenon~B as a talismanic
inscription in the Picatrix sense) is developed in
\S\ref{sec:verification}--\S\ref{sec:discussion}. The present section is
about the \emph{empirical} separation only.

\subsection{Phenomenon~A: the general categorial labeling system}
\label{sec:phenomenonA}

Phenomenon~A is the manuscript's pervasive positional categorial morphology
documented in \S\ref{sec:unified}. Its empirical signature is:

\begin{itemize}[itemsep=0pt,topsep=2pt]
\item canonical position~1 dominated by atom~A1 (with structurally similar
A2/A3/K1/Aa as minority alternatives);
\item canonical position~2 dominated by Q1/Q2 (the season-encoder slot of
nb71);
\item canonical positions~3--4 occupied by J-class and categorial atoms;
\item a closing A2 marker in roughly $80\%$ of canonical instances;
\item words drawn from the manuscript's general atom inventory of $\sim
166$ STA1 atom types.
\end{itemize}

This template is documented for the zodiac block (notebooks 67--71), and
\S\ref{sec:unified} extended it to the central nymph labels of f57v
(\texttt{@L0}), to the radial sun-ray labels of f70r1 (\texttt{@Ri}), to the
ring text of f70r2 (\texttt{@Cc}), and to body-text occurrences in the
herbal, pharmaceutical and recipe sections. Phenomenon~A therefore appears
to operate \emph{throughout most of the manuscript}, irrespective of section
or page genre, and is at minimum present in every textual element of f57v
\emph{except} the four concentric text rings.

\subsection{Phenomenon~B: the localized enumerative inscription}
\label{sec:phenomenonB}

Phenomenon~B is the textual layer documented in \S\S\ref{sec:subalphabet}
and \ref{sec:f70r}. Its empirical signature is the exact opposite of
Phenomenon~A on every dimension:

\begin{itemize}[itemsep=0pt,topsep=2pt]
\item it is built from a closed inventory of eight glyphs (the
\emph{private subalphabet} of Table~\ref{tab:subalphabet}), of which five
appear nowhere else in the corpus and three are concentrated on f57v at
enrichment factors of $75\times$--$302\times$;
\item it is structurally enumerative, not running text: in its densest
locus (line f57v.3) every ``word'' is a single isolated atom, $68/68 =
100\%$ isolated, the only such line in the manuscript;
\item it has cyclic regularity ($17 \times 4$ atomic characters in
f57v.3) consistent with a fixed-length inscription rather than free
discourse;
\item it does not follow Phenomenon~A's canonical positional template:
the subalphabet glyphs do not occupy zodiac p1/p2/p3/p4 slots and they do
not appear in any canonical Phenomenon~A label, including the
A1Q1J1A3C1 word of \S\ref{sec:shared_word};
\item geographically, it is confined to the four concentric text rings of
f57v (loci tagged \texttt{@Cc}/\texttt{+Cc}, lines f57v.2--f57v.5), and
to no other folio in the entire manuscript---confirmed by the
cosmological-suite test of \S\ref{sec:f70r}.
\end{itemize}

Phenomenon~B is therefore strictly localized, structurally enumerative,
and lexically disjoint from Phenomenon~A. Table~\ref{tab:two_phenomena}
summarises the contrast.

\begin{table}[h]
\centering
\small
\begin{tabular}{lll}
\toprule
Property & Phenomenon~A & Phenomenon~B \\
\midrule
Atom inventory & General STA1 ($\sim 166$ atoms) & Closed: 8 atoms (Tab.~\ref{tab:subalphabet}) \\
Geographic extent & Manuscript-wide & f57v concentric rings only \\
Structural form & Positional words (4--9 atoms) & Isolated single-atom tokens \\
Densest locus & ---  & f57v.3 ($17\!\times\!4$, $68/68$ isolated) \\
Canonical morphology (\S\ref{sec:zodiac_canon}) & Yes (zodiac canon) & No \\
Position-1 dominance & A1 ($85\%$ zodiac, $100\%$ \texttt{@L0}) & N/A \\
Test against cosmo suite (\S\ref{sec:f70r}) & Present in 14/14 folios & Absent in 14/14 folios \\
On f57v itself & \texttt{@L0}/\texttt{@Ro} labels & \texttt{@Cc}/\texttt{+Cc} ring text \\
Hypothesised function & Descriptive categorial labeling & Operative talismanic inscription \\
\bottomrule
\end{tabular}
\caption{Side-by-side empirical comparison of Phenomenon~A (the general
categorial labeling system) and Phenomenon~B (the localized enumerative
inscription on f57v). The two phenomena are disjoint on every measurable
dimension and coexist on f57v in spatially separated page elements: the
central nymph and radial labels carry Phenomenon~A, the four concentric
text rings carry Phenomenon~B.}
\label{tab:two_phenomena}
\end{table}

\subsection{Why the separation matters: f57v as the only co-presence}
\label{sec:why_separation}

The decisive empirical observation is that the two phenomena coexist
\emph{on the same folio without overlapping}. Folio f57v has thirteen
text loci (Table~\ref{tab:perline}), and they are partitioned exactly
into two groups by the categorisation of Table~\ref{tab:two_phenomena}:
the eight \texttt{@L0}/\texttt{@Ro} labels carry Phenomenon~A and zero
subalphabet atoms; the four \texttt{@Cc}/\texttt{+Cc} concentric ring
lines carry Phenomenon~B and zero canonical Phenomenon~A morphology.
The two layers occupy disjoint page elements.

This is the structurally non-trivial fact. The coexistence is not a
random mix of two systems but a \emph{spatial assignment}: descriptive
labeling around the figures, and a closed enumerative inscription on
the concentric rings around them. As the cosmological-suite test of
\S\ref{sec:f70r} confirms, f57v is the only folio in the entire
224-folio corpus where this pattern of co-presence occurs.

The argument that follows in
\S\S\ref{sec:verification}--\ref{sec:discussion} is that this exact
pattern---a descriptive astrological text accompanying a figure together
with a reserved talismanic inscription in non-standard characters
attached to a specific planetary position---is precisely what
\textit{Picatrix Latinus} prescribes for the manufacture of planetary
images, and that f57v is plausibly an instance of that practice. But the
\emph{empirical} separation of Phenomenon~A from Phenomenon~B does not
depend on the Picatrix interpretation. It is a fact about the manuscript
that has to be accounted for under any interpretive framework.


% ============================================================
\section{Direct verification against \textit{Picatrix Latinus} MS 793}
\label{sec:verification}

\S\S\ref{sec:subalphabet}--\ref{sec:two_phenomena} establish, on purely
internal evidence from the Voynich Manuscript, that folio f57v contains a
localized enumerative inscription (Phenomenon~B) coexisting on the same
folio with the manuscript's general categorial labeling system
(Phenomenon~A) and that this co-presence does not occur on any other folio
of the codex. The present section turns to external evidence and asks
whether any documented medieval textual tradition prescribes \emph{exactly
this kind of co-presence}---a descriptive astrological text accompanying a
figure together with a reserved talismanic inscription in non-standard
characters attached to a specific planetary position. The answer, we
argue, is yes: the tradition is the Latin Picatrix, the
\textit{karacteres} prescription specifically, and the strongest documentary
witness is the only fully illustrated extant Latin Picatrix manuscript,
Cracow Jagiellonian Library MS~793. We document four direct
correspondences with f57v, drawing on the photographic facsimile of MS~793
obtained in April~2026 and on the printed critical Latin edition of the
Picatrix text by \citet{pingree1986}.

\subsection{The MS~793 facsimile and its provenance}
\label{sec:facsimile}

The photographic facsimile of \bjms{} consulted here is a complete-codex
digital reproduction of $401$ pages, covering the entire bound
miscellaneous volume in which the Picatrix Latinus is the longest item.
The Picatrix section proper occupies approximately ff.~171r--197r of the
manuscript (= facsimile pp.~340--395), preceded by an unrelated
miscellaneous block of medical, astrological and computus material in
ff.~1--170 (cf.\ task~\#16 of the project repository, ``index of the
\bjms{} facsimile by manuscript folio'', currently in progress).

The manuscript itself is dated $1471$--$1474$ on internal grounds and is
attributed to Piotr Gaszowiec (1430--1474), court physician and
astrologer to King Casimir~IV Jagiellon, who copied the text in Cracow.
\citet{sniezynska2009}~(pp.~11--13) reconstructs that Gaszowiec did not
copy from a Cracovian original but from an earlier Bohemian prototype no
longer extant, dated c.~1390--1410 and produced in the workshop of King
Waclav~IV of Bohemia. The full transmission chain Gaszowiec~$\rightarrow$
Bohemian prototype~$\rightarrow$ Latin Picatrix~$\rightarrow$ Toledan
translation from the Castilian of Alfonso~X is discussed in
\S\ref{sec:transmission} below.

For the purposes of this section, the relevant facts are: (i)~MS~793 is
the only fully illustrated copy of the Latin Picatrix to survive; (ii)~it
is the closest extant witness in geography and date to the Voynich
Manuscript's C14 window (Voynich: 1404--1438; MS~793: 1471--1474, with a
prototype c.~1390--1410 that overlaps the lower bound of the Voynich
C14); (iii)~its illustrative program is partially abridged with respect
to the original Picatrix prescriptions, in ways that we will document
below.

\subsection{The Cracow Picatrix milieu (1410--1474)}
\label{sec:cracow_milieu}

The 1471--1474 date of MS~793 is not the date at which Picatrix entered
Cracow. \citet{sniezynska2009}~(pp.~12--13) documents continuous
Picatrix-related activity in the Cracow milieu from at least 1410:

\noindent (i)~A 1430 court document concerning Anna of Cilli (Anne, second
wife of Wladyslaw~II Jagiello) mentions a certain \textit{Nicolaus},
described as practising \textit{ars nigromantiae}, who reportedly
manufactured \textit{figure solis et lune} engraved on rock crystal---an
operative talismanic practice in the canonical Picatrix mould.

\noindent (ii)~Two Cracow practitioners, \textit{Henricus Bohemus} (Henryk
Czech) and Stanis\l aw of Kazimierz, are documented in the first half of
the fifteenth century as \textit{nigromantici} in court records.

\noindent (iii)~Jan of G\l og\'ow (1445--1507) wrote a
\textit{Tractatus de XLVIII imaginibus celestibus} that explicitly
discusses the manufacture of planetary talismans following Picatrix
prescriptions.

\noindent (iv)~Piotr Gaszowiec, the scribe of MS~793, was court
physician-astrologer to Casimir~IV Jagiellon and a professor at the
Cracow University faculty of medicine---a member of the same intellectual
milieu in which the Picatrix tradition was actively practised, not an
isolated antiquarian.

The Cracow Picatrix tradition is therefore documented \emph{anterior} to
MS~793 by approximately sixty years and predates the Voynich Manuscript's
lower C14 bound by only a few years. The Bohemian prototype that
Gaszowiec copied from (c.~1390--1410) overlaps the Voynich's C14 window
directly. This places the Cracow--Bohemia Picatrix milieu in the same
geographic and chronological neighbourhood as the Voynich's
manufacture, and provides the historical setting for the iconographic
correspondences documented in the next subsections.

\subsection{Visual confirmation~1: the central solar disc}
\label{sec:vc1_disc}

The most immediately obvious iconographic correspondence is the central
solar disc. Folio f57v of the Voynich shows, at the geometric centre of
its concentric ring system, a small but unambiguous schematic symbol: a
circle with a smaller circle inside and a dot in the centre---the
canonical astronomical symbol for the Sun, $\odot$.

On MS~793 f.~191r (= facsimile p.~382), the Picatrix text presents
``Forma Solis secundum opinionem'' in four versions corresponding to four
different astrological authorities. The right-hand figure on the page is
the Sun on his chariot (\textit{sol in curru}), drawn naturalistically as
an enthroned ruler. In his raised right hand he holds a golden disc with
a smaller dark circle inside and a central dot---naturalistically rendered
as a piece of golden jewellery, but iconographically the same object
$\odot$ that the Voynich draws schematically.

This is one of the few cases in the entire iconographic comparison where
the two manuscripts are not just \emph{topically} related but
\emph{literally} drawing the same object. The Voynich draws the symbolic
reduction; MS~793 draws the naturalistic version. The image files
\filename{historic/picatrix_p382_sun_disc.jpg} and
\filename{historic/f57v_diagram.jpg} document the side-by-side comparison
(see Figure~\ref{fig:disc}).

\begin{figure}[h]
\centering
\includegraphics[width=0.30\textwidth]{f57v_diagram.jpg}\hspace{1em}
\includegraphics[width=0.30\textwidth]{picatrix_p382_sun_disc.jpg}
\caption{Left: Voynich f57v, central solar disc $\odot$. Right: \bjms{}
f.~191r (facsimile p.~382), the right-hand Sun figure's hand-held disc.
The two are iconographically the same object: a circle within a circle
with a central dot, drawn schematically in the Voynich and
naturalistically in MS~793.}
\label{fig:disc}
\end{figure}

\subsection{Visual confirmation~2: the \textit{karacteres} on the Sun shields}
\label{sec:vc2_shields}

The same MS~793 folio (f.~191r, facsimile p.~382) contains two Sun
figures: the right-hand Sun in his chariot just discussed, and a
left-hand Sun figure on horseback. \emph{Both} figures hold shields, and
\emph{both} shields are inscribed with abstract compound sigils that are
neither Latin letters nor Greek letters nor any standard heraldic device.
The sigils are constructed from short straight strokes, small circles,
right-angle hooks and curved tails---the geometrical vocabulary of what
\citet{sniezynska2009}~(p.~103) explicitly describes as
\textit{karacteres}: ``transformations of the Arabic or Semitic alphabet,
encountered as medieval mason's marks'', noting (p.~105) that ``the use
of the magical alphabet and esoteric symbols is condemned in the
\textit{Speculum astronomiae}; perhaps this is why such inscriptions
contained in the text of \textit{Picatrix Latinus} appear in the
miniatures of the Cracow manuscript only once''.

That last observation is important: \citet{sniezynska2009} explicitly
notes that the Cracow scribe of 1471--1474 \emph{actively suppressed}
most of the karacteres prescribed in the Latin Picatrix text, in
compliance with the \textit{Speculum astronomiae}'s condemnation of
magical alphabets. The few karacteres that survived this scribal
censorship are precisely the ones we are now examining: the two compound
sigils on the Sun shields of f.~191r and the literal
$\mathrm{O\Lambda OIO\Lambda}$ inscription on the Venus comb of f.~191v.
What MS~793 preserves of the karacteres tradition is therefore a small
residue of what the original Picatrix prescribed.

The image files \filename{historic/picatrix_p382_left_shield.jpg} and
\filename{historic/picatrix_p382_right_shield.jpg} document the two
shield sigils in high resolution. They are reproduced side by side in
Figure~\ref{fig:shields}.

\begin{figure}[h]
\centering
\includegraphics[width=0.28\textwidth]{picatrix_p382_left_shield.jpg}\hspace{1em}
\includegraphics[width=0.28\textwidth]{picatrix_p382_right_shield.jpg}
\caption{The two compound karacter sigils on the Sun shields of \bjms{}
f.~191r (facsimile p.~382). Left: the shield of the left-hand Sun figure
(Sun on horseback). Right: the shield of the right-hand Sun figure (Sun in
chariot). Both sigils are constructed from short strokes, small circles
and angular hooks of a geometrically restricted vocabulary, identified by
\citet{sniezynska2009}~(p.~103) as \textit{karacteres} of the Picatrix
talismanic-alphabet tradition. They are two of only three surviving
karacteres in the entire MS~793 illustrative program, the third being
the $\mathrm{O\Lambda OIO\Lambda}$ inscription on the Venus comb of
f.~191v (\S\ref{sec:oloiol}).}
\label{fig:shields}
\end{figure}

\subsection{The literal $\mathrm{O\Lambda OIO\Lambda}$ inscription on the wooden comb of Venus}
\label{sec:oloiol}

On MS~793 f.~191v (= p.~383 of the facsimile PDF), the
``\textit{Forma Veneris secundum opinionem sapientis Picatricis}'' figure
holds a \textbf{wooden comb} (\textit{pectinem ligneum similem tabule})
inscribed with the magical word $\mathrm{O\Lambda OIO\Lambda}$. The object
is iconographically a comb, one of Venus's standard attributes alongside
the apple, with the inscription engraved on its flat central surface and
two sets of teeth visible above and below. Unlike the karacter sigils of
the Sun shields (\S\ref{sec:vc2_shields}), this inscription is preserved
\textbf{literally and verbatim} in the printed critical edition of the
Latin text by \citet{pingree1986}, where it appears in two distinct
passages of Book~II, Capitulum~X
(``\textit{De lapidibus appropriatis cuilibet planete et formacionibus
figurarum}''):

\medskip

\noindent\textit{Pingree 1986, p.~67, entry 27} (descriptive cataloguing of the
iconographic forms of Venus, Picatrix's own ``opinion''):

\begin{quote}
\itshape
``Forma Veneris secundum opinionem sapientis Picatricis est forma mulieris
tenentis in eius dextra manu malum et in sinistra \textbf{pectinem similem
tabule et istis figuris scriptum:} $\mathrm{O\Lambda OIO\Lambda}$.
Et hec est eius forma.''
\end{quote}

\noindent (``The form of Venus according to the opinion of the wise Picatrix is
the form of a woman holding in her right hand an apple and in her left hand
\textbf{a comb similar to a tablet, inscribed with these figures:}
$\mathrm{O\Lambda OIO\Lambda}$. And this is its form.'')

\medskip

\noindent\textit{Pingree 1986, p.~70, entry 55} (operative talismanic
instructions for the manufacture of a Venus love-amulet):

\begin{quote}
\itshape
``Ymago Veneris. Si ex formis Veneris feceris formam mulieris cuius corpus sit
humanum, caput vero avis necnon et pedes aquile, in dextra manu malum, in
sinistra vero \textbf{pectinem tenentis ligneum similem tabule talibus figuris
scriptum:} $\mathrm{O\Lambda OIO\Lambda}$, qui hanc ymaginem secum portaverit
bene recipietur et ab omnibus diligetur.''
\end{quote}

\noindent (``Image of Venus. If from the forms of Venus you make the figure of
a woman whose body is human but whose head is that of a bird and whose feet
are those of an eagle, holding an apple in her right hand and in her left
\textbf{a wooden comb similar to a tablet, inscribed with these figures:}
$\mathrm{O\Lambda OIO\Lambda}$, whoever carries this image with him will be
well received and loved by all.'')

\medskip

\noindent Three observations from the verbatim Latin text:

\noindent (i)~\textbf{The inscription is preserved as Greek capital
characters} in the printed critical edition. Pingree uses the unicode
characters $\mathrm{O\Lambda OIO\Lambda}$, not Latin letters that resemble
Greek. This is significant because it indicates that in the manuscripts
collated by Pingree (which include MS 793 among other Latin Picatrix
witnesses), the inscription was written either as actual Greek letters or as
signs that the editor judged equivalent to Greek letters --- in any case,
\textbf{distinct from the Latin script of the surrounding body of the text}.
This is exactly the category ``invented or transformed alphabet inserted as a
talismanic inscription'' that the Picatrix tradition prescribes.

\noindent (ii)~\textbf{The object is iconographically a wooden comb}, not a
tablet. Pingree's text twice writes \textit{pectinem similem tabule} (``a comb
similar to a tablet'', entry 27) and \textit{pectinem ligneum similem tabule}
(``a wooden comb similar to a tablet'', entry 55). The comb is one of the
standard iconographic attributes of Venus in medieval art, alongside the
apple, the mirror, and the rose. The flat surface of the comb between its two
sets of teeth carries the inscription. \textbf{This explains the visual
structure of the object in MS 793 f. 191v}: the user-observed ``8 vertical
strokes above and 12 vertical strokes below the central inscription'' (visible
in \filename{historic/picatrix_p383_tablet_only.jpg}) are the \textbf{teeth of
the wooden comb}, with 8 broad teeth on one side and 12 fine teeth on the
other in a typical medieval double-comb design.

\noindent (iii)~\textbf{The inscription has explicitly operative talismanic
function}. Entry 55 is not descriptive but \emph{prescriptive}: it is a recipe
for fabricating a Venus love-amulet, instructing the operator that a wooden
comb engraved with $\mathrm{O\Lambda OIO\Lambda}$ on its flat surface, when
worn, makes the bearer ``well received and loved by all'' (\textit{bene
recipietur et ab omnibus diligetur}). The thematic association is canonically
Venusian: the planet of love and attraction is invoked through an inscription
on her standard iconographic attribute (the comb). The Cracow scribe of MS
793, who suppressed nearly all other karacteres prescribed by the Picatrix
text in compliance with the \textit{Speculum astronomiae}'s condemnation of
magical alphabets \citep[see][p.~105]{sniezynska2009}, preserved this single
inscription verbatim --- presumably because of its central iconographic role
in the Venus images of the chapter.

\medskip

The two Picatrix passages together establish that, in the medieval
astrological-talismanic tradition into which this paper places the Voynich,
the use of inscribed non-Latin characters as operative talismanic devices
was \emph{textually documented and prescribed}, not merely incidental.
The f57v.3 alphabet ring of the Voynich ($17\times 4 = 68$ atomic characters
in cyclic structure) is a much more extensive instance of the same category
of textual phenomenon than the six characters of
$\mathrm{O\Lambda OIO\Lambda}$, but both belong to the same documented
tradition.

\subsection{Side-by-side comparison: alphabet versus compound sigils}
\label{sec:vc4_alphabet}

The four preceding subsections established three independent
correspondences between Voynich f57v and \bjms{} ff.~191r--191v: the
central solar disc (\S\ref{sec:vc1_disc}), the karacter sigils on the
Sun shields (\S\ref{sec:vc2_shields}), and the literal
$\mathrm{O\Lambda OIO\Lambda}$ inscription on the Venus comb
(\S\ref{sec:oloiol}). Each correspondence is documentary, anchored to a
specific page of MS~793 and to a specific passage of the Latin Picatrix
edited by \citet{pingree1986}. We now present a fourth correspondence
that is structural rather than iconographic: a direct visual comparison
of the f57v concentric ring text with the karacter sigils of MS~793,
designed to test what kind of textual relationship the two materials
enter into.

\subsubsection*{The polar unwrap of the f57v second ring}

The four concentric text rings of f57v are difficult to compare directly
with linear text because they are physically circular. To bring them into
a directly comparable form we developed an open-source polar
rectification tool (\filename{tools/circular_text_unwrap.ipynb} in the
project repository), which uses bilinear interpolation
(\texttt{scipy.ndimage.map\_coordinates}) to remap an annular region of a
circular diagram into a horizontal strip. The tool is fully parametric in
the centre coordinates and inner/outer radii of the annulus and is
provided as a standalone Jupyter notebook so that the experiment can be
reproduced or extended to other Voynich circular diagrams.

Applied to the second concentric text ring of f57v---the densest locus
of Phenomenon~B, line f57v.3 with its $17\!\times\!4$ structure of
single-atom isolated tokens---the polar unwrap produces a horizontal
strip in which the 68 atomic characters become individually legible
(Figure~\ref{fig:unwrap}). The image \filename{historic/f57v_ring2_alphabet_strip.jpg}
records the rectified strip; the four cyclic repetitions of the
17-character pattern are visible by inspection along the horizontal
axis.

\begin{figure}[h]
\centering
\includegraphics[width=0.95\textwidth]{f57v_ring2_alphabet_strip.jpg}
\caption{Polar unwrap of the second concentric text ring of Voynich
f57v (line f57v.3, the densest locus of Phenomenon~B), produced with
\texttt{circular\_text\_unwrap.ipynb}. The original circular
ring is remapped to a horizontal strip in which the 68 atomic
characters become legible as a sequence. Four cyclic repetitions of a
17-character pattern are present along the horizontal axis, matching
the $17\!\times\!4$ structure detected on purely textual grounds in
Table~\ref{tab:perline}.}
\label{fig:unwrap}
\end{figure}

\subsubsection*{Side-by-side: f57v.3 atomic characters vs MS~793 shield sigils}

We now compare the rectified f57v.3 strip with the two MS~793 shield
sigils of \S\ref{sec:vc2_shields}. We want to be explicit, before
presenting the comparison, about what kind of claim it supports. The
present subsection is a \emph{qualitative, side-by-side visual
comparison}, not a quantitative test. We attempted a reproducible
pixel-level quantification of the shared graphic vocabulary, including
scale-normalised binarisation, connected-component extraction, skeleton
decomposition of the compound sigils into elementary sub-strokes, and a
primitive-category distribution comparison via Jensen--Shannon
divergence with permutation testing (companion notebook~74b of the
project repository). The attempt proved not to be supportable on the
available source material: the photographic facsimile of MS~793 is
sufficiently noisy at the stroke-thickness scale that automated
segmentation cannot cleanly separate the karakter ink from the
parchment texture of the underlying shield paint, and the resulting
primitive distributions were dominated by segmentation artefacts rather
than by real stroke-shape statistics. We therefore present the
comparison below as a visual judgement, and retain the quantitative
test only as documented future work (cf.\ \S\ref{sec:future_work}) for
the higher-resolution archival imagery of MS~793 that would be obtained
by an in-person inspection at the Biblioteka Jagiello\'nska.

With that caveat stated, the comparison (Figure~\ref{fig:karacteres})
makes one structural property immediately clear to the eye: the visual
vocabulary used to construct individual characters of the Voynich ring
is, by inspection, the \emph{same vocabulary} from which the MS~793
shield sigils are built. Both inventories use short straight strokes,
small circles, right-angle hooks, and short curved tails as their
elementary graphic units; both avoid alphabetic letterforms and standard
heraldic devices; both place individual characters in spatial isolation
rather than ligated cursive sequences.

\begin{figure}[h]
\centering
\includegraphics[width=0.95\textwidth]{comparison_karacteres_v2.jpg}
\caption{Side-by-side comparison of the rectified f57v.3 ring strip
(top) with the two MS~793 Sun-shield sigils of f.~191r (bottom). The
elementary graphic vocabulary---short strokes, small circles, angular
hooks, curved tails---is shared between the two sources. The
relationship is one of \emph{atomic alphabet} (Voynich) to \emph{compound
sigils written in that alphabet} (MS~793), not of direct textual
copying.}
\label{fig:karacteres}
\end{figure}

\subsubsection*{Interpretation: alphabet vs compound sigils}

The structural relationship is asymmetric. The MS~793 shield sigils are
\emph{compound}: each shield carries a single sigil that is itself
constructed from several elementary marks combined into a unit. The
Voynich f57v.3 ring is \emph{atomic}: each of its 68 characters is a
single isolated mark, and the four cyclic repetitions show the
inscription as a fixed-length sequence of such atoms rather than as a
single compound device.

This is the structural pattern of an alphabet versus the words written
\emph{in} that alphabet. If the medieval karacteres tradition is taken
literally---an alphabet of geometric characters from which talismanic
inscriptions are composed---then the Voynich f57v.3 ring is a candidate
for the \emph{alphabet itself}, while the MS~793 shield sigils and the
$\mathrm{O\Lambda OIO\Lambda}$ comb inscription are candidates for
\emph{compound inscriptions} drawing on the same atomic vocabulary.

We do \emph{not} claim that the specific atomic characters of the f57v.3
ring can be matched one-to-one with the elementary components of the
MS~793 sigils. Such a one-to-one matching would require either a
key-document we do not possess or a much larger comparative corpus of
karacteres, and the present paper does not undertake it (cf.\
\S\ref{sec:limitations}). What we do claim, on the strength of the
\emph{visual} comparison in Figure~\ref{fig:karacteres}, is that the
two sources belong to the same \emph{graphic genre}, with the Voynich
ring carrying the more complete inventory and the MS~793 shields
carrying compound instances of that inventory. The claim is deliberately
calibrated to the evidence: it is a qualitative reading of the side-by-side
figure by a human observer, not a statistical test, and
\S\ref{sec:limitations}(a) lists it among the limitations that future
work at higher source resolution should revisit.

If this reading is correct, then the relationship of the Voynich to
\textit{Picatrix Latinus} is not one of textual descent but of
\emph{independent witnessing}: the Voynich preserves an entire alphabet
of karacteres on a single concentric ring, while MS~793 preserves only
three small relics of the same tradition (two compound shield sigils and
one literal six-character word) after its scribe's
\textit{Speculum astronomiae}-driven suppression of the rest. The
implications of this asymmetric preservation pattern are developed in
\S\ref{sec:more_complete} below.


% ============================================================
\section{Discussion: the Voynich as an independent witness of the
\textit{Picatrix} tradition}
\label{sec:discussion}

The empirical findings of \S\S\ref{sec:subalphabet}--\ref{sec:two_phenomena}
and the iconographic correspondences of \S\ref{sec:verification} are
independent of any single interpretive frame and stand as facts about
the manuscript whatever their ultimate explanation. This section
proposes the interpretive frame we believe best accounts for them, and
explicitly states what the frame does \emph{not} claim. The proposal
has four moving parts: (i)~a methodological reformulation of the
``Picatrix hypothesis'' from textual derivation to independent
witnessing; (ii)~the observation that the Voynich preserves the
\textit{karacteres} tradition more completely than its own closest
documented contemporary; (iii)~the placement of the Voynich within
the documented Toledo--Prague--Cracow transmission chain; and
(iv)~an honest enumeration of the limitations the present paper does
not resolve.

\subsection{Reformulation: independent witness, not derivation}
\label{sec:reformulation}

The traditional ``Picatrix hypothesis'' for the Voynich, going back to
\citet{dimperio1978}, has typically been formulated as a textual
descent claim: \emph{the Voynich is, in some sense, a copy or
translation of the Latin Picatrix or of one of its sources}. So
formulated, the hypothesis is methodologically weak and easily
rejectable. The two companion papers of this series~\citep{jimenez2026negative,
jimenez2026twotier} have already established that no substitution-based
mapping of Voynich glyphs onto any natural language tested produces
readable text against any of the obvious source candidates, and that
the Voynich's positional structure is consistent with a designed
combinatorial code rather than with the surface form of an underlying
natural-language source. A direct textual descent of the Voynich from
the Picatrix Latinus is not compatible with these findings.

We therefore reformulate the hypothesis. The Voynich is \emph{not}
proposed as a textual derivation of Picatrix Latinus or of any
specific Picatrix manuscript. It is proposed as an \emph{independent
witness} of the same broader medieval astrological-talismanic tradition
that the Picatrix corpus codifies. Specifically:

\begin{itemize}[itemsep=0pt,topsep=2pt]
\item The Picatrix text (in Arabic and Latin recensions) is one
documented codification of an older tradition of categorial astrological
texts plus reserved talismanic inscriptions in non-standard characters
(\textit{karacteres}).
\item Cracow Jagiellonian Library MS~793 is one documented illustrated
witness of the Latin recension of that codification, datable to
1471--1474 with a Bohemian prototype c.~1390--1410.
\item The Voynich Manuscript, on the textual evidence of
\S\S\ref{sec:subalphabet}--\ref{sec:two_phenomena}, contains a
\emph{co-presence} of the same two textual layers (descriptive
categorial text plus a reserved enumerative inscription in
non-standard characters) on a folio whose iconography
(\S\ref{sec:verification}) overlaps directly with the most karacteres-rich
folios of MS~793.
\end{itemize}

The reformulated hypothesis is therefore a claim of \emph{shared
membership in a documented tradition}, not of textual lineage.
Crucially, this kind of claim is much weaker and much easier to
defend: it does not require us to identify a specific source manuscript,
to reconstruct a specific copying chain, or to demonstrate one-to-one
character correspondences. It requires only that the Voynich's
structural fingerprint be incompatible with random scribal invention
and compatible with the documented practice of the tradition---both
of which are empirical claims that the previous sections of this paper
defend.

\paragraph{A methodological note on ``extractable content''.} A common
reflex when evaluating work on the Voynich is to measure progress in
terms of how much readable text a proposal produces. Under that
measure, the present paper does not advance the field: no specific
message is recovered, no translation is offered, no continuous
passage is made legible. We want to be explicit that this absence is
\emph{not a residual gap} of the present analysis but an
\emph{architectural prediction} of the interpretive frame of
\S\ref{sec:verification}--\S\ref{sec:reformulation}. The medieval
\textit{karacteres} tradition, as documented by
\citet{sniezynska2009}~(p.~103) and by the prescriptive passages of
\citet{pingree1986}, is \emph{operative} rather than communicative:
the characters are deployed for planetary correspondence in a
talismanic device, not to transmit a linguistic message to a reader.
A successful identification of a folio with that tradition should
therefore predict that no linguistic decoding will recover content,
which is exactly what the nineteen substitution-based models of
\citet{jimenez2026negative} find. The two results are mutually
coherent: the ``failure to read'' of the first paper and the
``operative inscription'' of the present paper are the same fact seen
from two angles. We do not, therefore, present the silence of
specific content as a limitation of the interpretive proposal; we
present it as one of its predictions, and as the reason that
decipherment-shaped evaluations are the wrong test for the reading
developed here.

\subsection{The Voynich as a more complete witness than MS~793}
\label{sec:more_complete}

A striking consequence of the reformulation in \S\ref{sec:reformulation}
is that, on the karacteres dimension specifically, the Voynich is
arguably a \emph{more complete} witness of the medieval talismanic
tradition than MS~793 itself. As documented in \S\ref{sec:vc2_shields}
and \S\ref{sec:oloiol}, the entire surviving karacteres material in
MS~793's illustrative program consists of three vestiges: the two
compound Sun-shield sigils of f.~191r and the literal
$\mathrm{O\Lambda OIO\Lambda}$ inscription on the Venus comb of
f.~191v. The Voynich, by contrast, carries on a single folio (f57v,
lines f57v.2--f57v.5) an entire eight-glyph private subalphabet
deployed in a structurally enumerative $17\!\times\!4$ inscription.
By counts, this is more karacteres material on a single folio than
MS~793 preserves in its entire illustrative program.

This is not a claim of priority or of greater authenticity, but of
asymmetric preservation: the same tradition shows up in the Voynich as
a relatively complete inscription and in MS~793 as three small relics.
We emphasise that this asymmetry is offered as a \emph{complementary
historical hypothesis}, not as independent evidence for the
reformulated reading of \S\ref{sec:reformulation}. The load-bearing
evidence of this paper is the $\sim 10^{-26}$/$\sim 10^{-64}$
concentration of the private subalphabet (\S\ref{sec:probability}) and
the four positive iconographic correspondences of
\S\ref{sec:verification}; the asymmetry discussed here is a
downstream observation that, if accepted, would help rationalise why
so little karacter material survives in the Cracow witness, but the
paper's argument does not depend on it.

This caveat is important because the asymmetry claim is partly an
\textit{argumentum ex silentio}: it reasons from what MS~793 does
\emph{not} show in order to support a reading of what the Voynich
\emph{does} show. \citet{sniezynska2009}~(p.~105) documents that the
Cracow scribe actively suppressed most of the karacteres prescribed by
his Latin source text in compliance with the \textit{Speculum
astronomiae}'s condemnation of magical alphabets, which supplies a
concrete and historically situated mechanism for the gap we observe in
MS~793; but gap-based arguments are in general weaker than
presence-based arguments, and a stricter reader is entitled to
discount this particular line and still accept the main empirical
claim. The temporal direction of the asymmetry, under which MS~793
was copied several decades after the Voynich's C14 window and therefore
under stronger scholastic pressure than the Voynich's manufacturer
plausibly faced, is offered here as the most natural historical
reading of that gap, not as a demonstrated fact.

\subsection{The Toledo--Prague--Cracow transmission chain}
\label{sec:transmission}

The reformulated hypothesis becomes historically concrete when placed
against the documented transmission chain of astrological manuscripts
in Central Europe between the late thirteenth and the late fifteenth
century. \citet{sniezynska2009}~(\textit{Wprowadzenie}, pp.~11--13)
documents the chain in some detail; we summarise the four nodes most
relevant to the Voynich's C14 window.

\noindent\textbf{Toledo (mid-13th century).} The Latin translation of
\textit{Gh\=ayat al-\d{H}ak\={\i}m} as \textit{Picatrix} is produced in
Castile in or around 1256, in the orbit of Alfonso~X's translation
workshop, from a Castilian intermediate translation. Alfonso~X is the
documented source from whom several Central European royal courts
acquired astrological manuscripts in the second half of the
thirteenth century.

\noindent\textbf{Prague (late 13th--early 15th century).} \v{P}\v{r}emysl
Otakar~II of Bohemia is documented as receiving astrological
manuscripts \emph{directly from Alfonso~X}. V\'aclav~II's court
astrologers, including Alvaro de Oviedo (d.~1299), reached Prague
``mainly from Toledo''~\citep[p.~12]{sniezynska2009}. Under Charles~IV
and his son V\'aclav~IV, Prague became the major Central European
centre of astrological manuscript production: Havel of Strahov,
Niccol\`o Beccari, and K\v{r}i\v{s}\v{t}an of Prachatice were active
in V\'aclav~IV's circle. The lost Bohemian prototype of MS~793,
reconstructed by Sniezynska-Stolot to c.~1390--1410, belongs to this
Prague phase.

\noindent\textbf{Cracow (1410--1474).} As documented in
\S\ref{sec:cracow_milieu}, the Picatrix tradition arrives in Cracow at
the latest by 1410 (Anna of Cilli's \textit{Nicolaus}) and is
continuously practised there through the fifteenth century, culminating
in Gaszowiec's MS~793 of 1471--1474.

\noindent\textbf{The Voynich's C14 window (1404--1438).} The Voynich's
radiocarbon dating overlaps the Bohemian phase of the chain almost
exactly. The Bohemian prototype that Gaszowiec later copied
(c.~1390--1410) is essentially contemporary with the Voynich. We do
not claim that the Voynich is that prototype, or a copy of it, or any
other specific manuscript. What we do claim is that the Voynich's
manufacture date sits squarely inside a documented Central European
geography and chronology in which the Picatrix tradition was actively
copied, illustrated, and operatively practised.

The transmission chain therefore provides a plausible historical home
for the Voynich. Under the reformulated hypothesis of
\S\ref{sec:reformulation}, the Voynich is one more witness of the same
broad textual tradition that flowed from Toledo through Prague to
Cracow---neither a copy of any specific surviving manuscript, nor an
isolated invention, but a sibling production of the same milieu.

\subsection{f57v within the wider cosmological suite}
\label{sec:cosmological_program}

A final implication of the reformulated hypothesis concerns the placement
of f57v within the wider cosmological/astronomical block of the Voynich.
The robustness test of \S\ref{sec:f70r} already established that f57v is
the unique locus of Phenomenon~B in that block: among the fourteen other
foldout panels of ff.~67--70, the seven strongest subalphabet glyphs do
not appear at all. This is consistent with the reading developed in
\S\S\ref{sec:reformulation}--\ref{sec:more_complete}: f57v is not an
alien folio bound into the codex by accident, but a member of the same
cosmological suite that nonetheless carries one additional textual layer
that no other folio in the suite carries.

The specific identification of the four central nymphs of f57v
(planets, seasons, cardinal directions, classical elements, or some
combination), and the corresponding identification of the four concentric
rings, is left as a separate question. The structural test of
\S\ref{sec:test_l0} shows that the four \texttt{@L0} labels follow the
canonical zodiac positional morphology and that position~5 exhibits a
$2$--$2$ binary split (A3/B1) consistent with a binary categorical
attribute, but the resolution of which binary attribute this is (e.g.\
solar/lunar, hot/cold, masculine/feminine, fertile/sterile) requires
external evidence beyond the present paper. The textual finding stands
independently of any such identification.

\subsection{Limitations and what this paper does not claim}
\label{sec:limitations}

The reformulated hypothesis is deliberately weaker than the
``decipherment'' claims that have historically been associated with
the Voynich Manuscript, and it is important to be explicit about
what this paper does \emph{not} claim. The empirical findings of
\S\S\ref{sec:subalphabet}--\ref{sec:two_phenomena} are robust to all
of the limitations below; only the interpretive frame of
\S\ref{sec:reformulation}--\S\ref{sec:transmission} is conditional on
external evidence that future work would have to settle.

\noindent\textbf{(a)~No one-to-one identification of glyphs, and no
quantitative vocabulary test.} We do not identify specific Voynich
subalphabet glyphs with specific characters of the MS~793 karacter
inventory. The graphical comparison of Figure~\ref{fig:karacteres} is
restricted to membership in a shared geometric vocabulary (short
strokes, small circles, hooks, curved tails) and does not assert that,
for example, a given Voynich atom corresponds to a given component of
the Sun-shield sigils. A defensible one-to-one matching would require
either a key document we do not possess or a substantially larger
corpus of medieval karacteres to provide statistical leverage.
Furthermore, the shared-vocabulary claim of \S\ref{sec:vc4_alphabet}
is itself offered as a qualitative visual reading of the side-by-side
figure, not as a quantitative statistical test. As reported in
\S\ref{sec:future_work}(d), we attempted to replace the visual
judgement with a reproducible pixel-level quantification of the
primitive-category distributions of the two sources, but the
available photographic facsimile of MS~793 is too noisy at the
stroke-thickness scale to support automatic segmentation, and the
companion notebook of that attempt is retained in the repository as
documentation of the failure mode. A stricter quantitative test
awaits higher-resolution archival imagery of MS~793 and a matched
external control source.

\noindent\textbf{(b)~No claim of textual derivation or authorship.} We
do not claim that the Voynich is a copy, translation, abridgement, or
derivative of any specific Picatrix manuscript or of any other
identified text. We do not claim a specific author, scribe, or place
of manufacture.

\noindent\textbf{(c)~No claim of underlying natural language.} We do
not claim that the Voynich, taken as a whole, encodes any specific
natural language. The first companion paper~\citep{jimenez2026negative}
established that nineteen substitution-based decipherment models
against five Semitic languages all converge on the same structural
ceiling without producing readable text, and the second companion
paper~\citep{jimenez2026twotier} established a positional code
structure inconsistent with simple natural-language transcription.
The interpretation of Phenomenon~A as ``categorial labeling'' and of
Phenomenon~B as ``operative talismanic inscription'' is compatible
with both findings.

\noindent\textbf{(d)~The Cracow milieu hypothesis is plausible, not
unique.} The Toledo--Prague--Cracow chain of \S\ref{sec:transmission}
is one documented pathway through which the relevant tradition was
demonstrably active in Central Europe during the Voynich's C14 window,
but it is not the only conceivable pathway. Other Central European
courts (Buda, Vienna, the imperial chancery, the German university
towns) had access to overlapping astrological material in the same
period, and the present paper does not exclude any of them.

\noindent\textbf{(e)~The interpretive frame is independent of the
empirical findings.} Most importantly, none of the empirical results
of \S\S\ref{sec:subalphabet}--\ref{sec:two_phenomena} depend on the
Picatrix interpretation being correct. The private subalphabet is a
fact about the STA1 corpus, the cosmological-suite test is a fact
about the IVTFF transcription, and the canonical-morphology compliance
of the f57v nymph labels is a fact about the positional analysis of
notebook~71. If the Picatrix interpretation of \S\ref{sec:reformulation}
is later refuted, the empirical findings remain in place and have to
be accounted for under whatever alternative frame is proposed.


% ============================================================
\section{Conclusions}
\label{sec:conclusions}

This paper has reported a single quantitative finding about the Voynich
Manuscript and has developed an interpretive proposal around it. We
recapitulate both, in the order in which they were established, and
then list the directions in which the work can most usefully be
extended.

\subsection{Summary of empirical findings}
\label{sec:summary_empirical}

The empirical core of the paper is that folio f57v of the Voynich
Manuscript carries a private subalphabet of eight STA1 atoms
(Table~\ref{tab:subalphabet}), of which five are exclusive to that
single folio of the entire corpus and three are concentrated on it
at enrichment factors of $75\times$--$302\times$ over baseline. The
joint probability of the concentration under the simplest null is
$\sim 10^{-64}$ (\S\ref{sec:probability}). The subalphabet lives
exclusively in the four concentric text rings of f57v, most densely
in line f57v.3 (the only line in the manuscript with $100\%$
isolated single-character ``words'', \S\ref{sec:intrafolio}), and is
absent from the eight nymph labels of the same folio. f57v is the
rank-$1$ folio of the entire manuscript by proportion of isolated
atoms, by a wide margin (\S\ref{sec:ranking}); and the
cosmological-suite robustness test of \S\ref{sec:f70r} confirms that
the anomaly is folio-specific rather than genre-specific.

In parallel, the four central nymph labels of f57v conform $4/4$ to
the canonical zodiac positional morphology of the author's previous
unpublished work (\S\ref{sec:test_l0}), and the recurring word that
those labels share with the zodiac block, the cosmological foldouts
and several body-text loci of the herbal and recipe sections
(\S\ref{sec:shared_word}) shows that this categorial labeling system
operates manuscript-wide. f57v therefore carries two textually
disjoint layers (Table~\ref{tab:two_phenomena}): a Phenomenon~A
labeling layer that is shared with the rest of the manuscript, and a
Phenomenon~B alphabet-ring inscription that is unique to the folio.
f57v is the only folio in the corpus where the two layers coexist.

\subsection{Summary of the interpretive proposal}
\label{sec:summary_interpretive}

The four iconographic correspondences against Cracow Jagiellonian
Library MS~793 documented in \S\ref{sec:verification} (the central
solar disc, the karacter sigils on the Sun shields, the literal
$\mathrm{O\Lambda OIO\Lambda}$ inscription on the Venus comb verified
verbatim against \citet{pingree1986}, and the structural
alphabet-vs-compound-sigils relationship revealed by the polar
unwrap) place Phenomenon~B within the documented \textit{karacteres}
prescription of the medieval astrological-talismanic tradition
codified by the Latin Picatrix.

The interpretive proposal of \S\ref{sec:reformulation} is that the
Voynich should be read as an \emph{independent witness} of that
tradition, not as a textual derivation of any specific Picatrix
manuscript. The proposal places the Voynich within the documented
Toledo--Prague--Cracow transmission chain of astrological manuscripts
in Central Europe between the late thirteenth and the late fifteenth
century (\S\ref{sec:transmission}); the Voynich's C14 window
($1404$--$1438$) overlaps the Bohemian phase of that chain almost
exactly. It is explicitly weaker than ``decipherment'' and is
independent of the empirical findings, which stand under any
alternative interpretive frame.

\subsection{Directions for future work}
\label{sec:future_work}

The empirical findings of this paper open several directions in which
the present analysis can be extended or independently verified.

\noindent\textbf{(a)~In-person archival inspection of MS~793.} The
present comparison relies on a complete-codex photographic facsimile
of MS~793 obtained in April~2026. An in-person inspection of the
original at the Biblioteka Jagiello\'nska in Cracow would allow direct
verification of the karacter sigils on the Sun shields and of the
Venus comb inscription, including any pigmentation, palaeographic
detail, or marginal annotations not visible in the digital facsimile.
Task~\#16 of the project repository (``index the full MS~793 facsimile
by manuscript folio'') is a prerequisite step toward this inspection.

\noindent\textbf{(b)~Comparison with the other ${\sim}20$ extant Latin
Picatrix manuscripts.} \citet{pingree1986} catalogues approximately
twenty surviving Latin Picatrix manuscripts in his census, of which
MS~793 is the only fully illustrated copy. Several of the others
contain partial illustrative material or marginal sigils that have not
been systematically compared with the Voynich. A targeted survey of
karacter inscriptions across the Pingree census would establish
whether the f57v subalphabet shares specific atoms with any extant
Picatrix witness, and would test the asymmetric-preservation claim of
\S\ref{sec:more_complete} against a wider sample.

\noindent\textbf{(c)~Palaeographic analysis of the Voynich script
against contemporary Bohemian and Cracow manuscripts.} The Voynich's
glyph forms have been compared against many Latin and vernacular
hands of the late fourteenth and early fifteenth century, but to the
present author's knowledge they have not been systematically compared
against the karacter inventories of contemporary Central European
talismanic manuscripts. The structural similarity documented in
Figure~\ref{fig:karacteres} suggests that such a comparison would be
diagnostically informative.

\noindent\textbf{(d)~Quantitative characterisation of the graphic
vocabulary shared between f57v and MS~793.} The side-by-side comparison
of Figure~\ref{fig:karacteres} establishes, at the qualitative level,
that the rectified f57v.3 ring and the MS~793 Sun-shield sigils draw
from a common inventory of elementary marks (short strokes, small
circles, right-angle hooks, curved tails), but the comparison is a
visual judgement and (as discussed in \S\ref{sec:vc4_alphabet}) a
sceptical reader is entitled to ask for quantitative support. We did
attempt such a quantification as a companion to this paper (companion
notebook~74b of the project repository, \texttt{74b\_f57v\_graphic\_%
vocabulary\_quantitative.ipynb}). The pipeline rescales all source
images to a common height, binarises with Otsu thresholding, extracts
connected components, decomposes the compound MS~793 sigils via
skeletonisation into elementary sub-strokes at skeleton junction
points, classifies each resulting stroke into one of six primitive
categories (dot, loop, straight stroke, arc, hook, complex), and
compares the f57v and shield primitive-frequency distributions with
Jensen--Shannon divergence under a permutation null. The attempt
failed at the image-quality step: the publicly available photographic
facsimile of MS~793 at its current resolution is too noisy, at the
stroke-thickness scale, for automatic segmentation to cleanly separate
karakter ink from the parchment texture of the underlying painted
shield. The resulting primitive distributions were dominated by
segmentation artefacts rather than by real stroke-shape statistics,
and the reported Jensen--Shannon divergence varied by a factor of two
under modest changes in the binarisation threshold. The notebook is
nonetheless retained in the repository as a documented methodological
attempt, both so that the procedure can be reproduced by any reader
and so that the exact failure mode is on the record. A successful
quantitative test of the shared-vocabulary claim therefore requires
(i)~higher-resolution archival imagery of MS~793 than is presently
available to the project, most plausibly obtained by an in-person
inspection at the Biblioteka Jagiello\'nska (see~(a) above); and
(ii)~a matched external control source (e.g.\ abstract marks from
unrelated late-medieval manuscripts, mason's marks from the same
geographic region) to establish a baseline against which the
f57v--MS~793 distance can be judged small or large. Until both of
those are in hand, the shared-vocabulary claim of
\S\ref{sec:vc4_alphabet} stands as a qualitative reading of the
side-by-side figure, not as a quantitative result.

\noindent\textbf{(e)~Textual identification of specific characters in
the f57v alphabet ring.} Beyond the vocabulary-level test of~(d), a
full character-level matching between the $17\!\times\!4$ characters
of line f57v.3 and the elementary components of the MS~793 shield
sigils, the $\mathrm{O\Lambda OIO\Lambda}$ comb inscription, and
karacter inventories in other Picatrix witnesses would test the
reformulated hypothesis at the level of individual glyph identity.
This is the most direct but also the most ambitious extension of the
present work, and it presupposes either a key-document we do not
possess or the successful construction of a substantially larger
comparative karacteres corpus from surviving Picatrix witnesses.

\noindent\textbf{(f)~Identification of the four central nymphs of
f57v.} The present paper has deliberately set aside the question of
which entities the four central nymphs of f57v represent (planets,
seasons, cardinal directions, classical elements, or some
combination). The structural analysis of \S\S\ref{sec:test_l0}
shows that the four \texttt{@L0} labels follow the canonical zodiac
positional morphology and that position-$5$ exhibits a
$2$--$2$ binary split (A3/B1) consistent with a binary categorical
attribute. A targeted analysis of the four labels against
candidate identification schemes (Sun--Venus--Mercury--Moon for the
inner planets; spring--summer--autumn--winter for the seasons;
etc.) is the natural follow-up.

\noindent\textbf{(g)~Application of the cosmological-suite test to
other circular diagrams of the Voynich.} The polar unwrap tool
\filename{tools/circular_text_unwrap.ipynb} of
\S\ref{sec:vc4_alphabet} is parametric in centre coordinates and
inner/outer radii and is therefore applicable to any Voynich
circular diagram. Systematic application of the same enumerative-vs-
running-text metric used in \S\ref{sec:intrafolio}, combined with
the rectified-strip visualisation, to all the manuscript's
concentric-text folios would test whether any other folio shows a
weaker version of the f57v signature that the present blunt
folio-level analysis has missed.

\subsection{Closing remark}

The Voynich Manuscript has been a difficult object for computational
methods because almost every published claim about it has been a
claim about the entire codex at once. The strategy of the present
paper has been the opposite: to find the single most distinctive
point in the manuscript and to ask what it points to. The point
turned out to be folio f57v, the eight glyphs of
Table~\ref{tab:subalphabet}, and the $17\!\times\!4$ enumeration of
line f57v.3; what it points to, on the documentary evidence of
\S\ref{sec:verification} and on the historical placement of
\S\ref{sec:transmission}, is a documented medieval textual category
in which the Voynich now appears to participate. The companion
papers~\citep{jimenez2026negative,jimenez2026twotier} together with
the present work suggest, on cumulative empirical and documentary
grounds, that the Voynich is best read as a medieval categorial
pasigraphy in the broad tradition codified by the Latin Picatrix---an
encyclopaedic categorial code carrying, on at least one of its
folios, an operative talismanic inscription of the same tradition's
\textit{karacteres}.


% ============================================================
\section*{Data and code availability}
\label{sec:data_availability}

The Voynich text used throughout this paper is the STA1~2.0 super
transliteration by Ren\'e Zandbergen~\citep{zandbergen_sta1}, available
at \url{https://voynich.nu/}, together with the IVTFF format
specification (Version~2.0.1, July~2025)~\citep{zandbergen_ivtff}. All
quantitative results, intermediate JSON outputs, and figures of this
paper are reproducible from the project's public GitHub repository,
which contains:

\begin{itemize}[itemsep=0pt,topsep=2pt]
\item \filename{74_f57v_private_subalphabet.ipynb} --- the principal
companion notebook implementing the per-folio glyph counting, the
probability calculation, the per-line breakdown of f57v, the
rank-ordering of the $224$-folio corpus, and the canonical-morphology
compliance tests of \S\S\ref{sec:test_l0}--\ref{sec:test_ri};
\item \filename{f57v_subalphabet_resultados.json} --- the structured
JSON output of the notebook, used as the source for
Tables~\ref{tab:subalphabet}, \ref{tab:perline}, \ref{tab:ranking},
\ref{tab:l0}, \ref{tab:ri}, and \ref{tab:shared_word};
\item \filename{f57v_subalphabet_cosmological_robustness.json} ---
the structured JSON output of the cosmological-suite robustness test
of \S\ref{sec:f70r}, used as the source for Table~\ref{tab:cosmotest};
\item \filename{tools/circular_text_unwrap.ipynb} --- the standalone
polar rectification tool of \S\ref{sec:vc4_alphabet}, parametric in
centre coordinates and inner/outer radii of the annulus, applicable to
any Voynich circular diagram;
\item the \filename{historic/} directory containing high-resolution
crops of the relevant pages of the MS~793 photographic facsimile
(\filename{picatrix_p382_*.jpg}, \filename{picatrix_p383_*.jpg}), the
high-resolution crops of Voynich f57v
(\filename{f57v_diagram.jpg}, \filename{f57v_ring2_alphabet_strip.jpg}),
and the side-by-side comparison images
(\filename{comparison_karacteres_v2.jpg}).
\end{itemize}

The Sniezynska-Stolot~2009 monograph is available open-access at the
Jagiellonian University institutional repository (RUJ). The Pingree~1986
critical edition of the Latin Picatrix is available at major research
libraries; the cited passages on \textit{Forma Veneris} (Book~II,
Capitulum~X, entries~$27$ and~$55$, pp.~$67$ and~$70$) were verified
directly against the printed edition.


% ============================================================
\section*{Acknowledgements}

The author thanks Ren\'e Zandbergen for the STA1 corpus and for
maintaining the IVTFF format specification, both of which made the
present analysis possible; Ewa \'Snie\.zy\'nska-Stolot, whose 2009
monograph on the iconography of \textit{Picatrix Latinus} in MS~793
provides the foundational scholarly source for the iconographic
comparison developed here; Mar\'ia Dolores Mart\'inez Gavil\'an
(Universidad de Le\'on, retired), whose 2025 article on Juan
Caramuel's \textit{Apparatus philosophicus} provides the theoretical
framework of \S\ref{sec:framework} and without which the present
paper would not have known how to name what it was seeing in the
corpus; the Beinecke Rare Book and Manuscript Library of Yale
University for making high-resolution images of MS~408 freely
available in the public domain; and the Biblioteka Jagiello\'nska in
Cracow for preserving MS~793 and for making its photographic
facsimile accessible. The author also acknowledges Anthropic Claude
(Opus~4.6, 1M~context) for tool-assisted code development, manuscript
drafting, and structured output generation; all content has been
verified by the author and all empirical claims are reproducible from
the public companion artefacts listed in
\S\ref{sec:data_availability}.

% End of body skeleton.
% Next step: full drafting section by section, in the recommended order:
% Section 3 → 4 → 5 → 6 → 7 → 8 → 2 → 1 → 9
